\section{Multi-Party Key Exchange in the Bounded Storage Model}
\label{sec:bsm-nike}

In this section, we extend the two-party NIKE in the BSM \cite{C:CacMau97} to a multi-party one and show the optimality of our construction up to constant factors within the natural class of schemes satisfying extendability and local key independence by giving a lower bound for the users' storage consumption. All algorithms in this section are considered as streaming word-RAMs except for the storage function applied by the adversary, which can be any unbounded function with limited output size.

\subsection{Our Construction}
%
%\subsection{Background of Bounded Storage Model}
%\label{sec:bsmke-bkg}
%\subsection{Our Construction}
%\label{sec:bsmke-con}

%
%

%We first refer the reader to \cref{sec:pairwisedef} for the definitions of uniformly and approximately pairwise independence, (approximately) strongly 2-universal hash function family, and 2-universal hash function family, and recall their instantiations in \cref{fig:con-hash}.~\footnote{The approximate pairwise independence and approximate 2-universal hash function family are defined in \cite{FOCS:Zuckerman91,C:CacMau97} in an implicit way.}


%
\heading{Pairwise independence.} We first recall the definitions of uniformly and approximately pairwise independence, (approximately) strongly 2-universal hash function family, and 2-universal hash function family, and recall their instantiations in \cref{fig:con-hash}.~\footnote{The approximate pairwise independence and approximate 2-universal hash function family are defined in \cite{FOCS:Zuckerman91,C:CacMau97} in an implicit way.}
%\section{Pairwise Independence}
%\label{sec:pairwise}
%We now recall the definitions of (approximate) pairwise independence, (approximately) strongly 2-universal hash function family, and 2-universal hash function family.
%
%\begin{definition}[Uniform pairwise independence]
%\label{def:upairwise}
%A sequence of random variables $x_1,\cdots,x_q$ with alphabet $\mathcal Y$ are uniformly pairwise independent if for any $1\leq i<j\leq q$ and any $\hat y_1, \hat y_2 \in \mathcal X$, we have
%	$$\Pr[x_i = \hat y_1 \land x_j = \hat y_2] = \frac{1}{|\mathcal{Y}|^2}.$$
%\end{definition}
%
%
%\begin{definition}[Approximate pairwise independence \cite{FOCS:Zuckerman91,C:CacMau97}]
%\label{def:apairwise}
%A sequence of random variables $x_1,\cdots,x_q$ with alphabet $\mathcal Y$ are approximate pairwise independent if for any $1\leq i<j\leq q$ and any distinct $\hat y_1, \hat y_2 \in \mathcal Y$, we have
%$$\Pr[x_i = \hat y_1 \land x_j = \hat y_2] = \frac{1}{|\mathcal{Y}|(|\mathcal{Y}|-1)}.$$
%%The definition of \emph{approximate pairwise independence} is exactly the same as \cref{def:upairwise} except that we replace 
%%``$\frac{1}{|\mathcal{Y}|^2}$''
%%by
%%``$\frac{1}{|\mathcal{Y}|(|\mathcal{Y}|-1)}$''.
%\end{definition}
%
%
%%For any $q\in\N$, a sequence of random variables $\vec x_1,\cdots,\vec x_q$ from a range $\mathcal X$ are pairwise independent if for any $1\leq i<j\leq q$ and any $\hy_1, \hy_2 \in \bin^\nn$ with $\hy_1\neq \hy_2$, we have
%%	$$\Pr[\vec x_i = \hy_1 \land \vec x_j = \hy_2] = \frac{1}{|\mathcal{X}|^2}.$$
%
%\begin{definition}[Strongly 2-universal hash function family]
%\label{def:shash}
%	A set $\mathcal H$ of functions with domain $\mathcal X$ and range $\mathcal Y$ is said to be a \emph{strongly 2-universal hash function family} if for any distinct $x_1,x_2 \in \mathcal{X}$, $h(x_1)$ and $h(x_2)$ are uniformly pairwise independent for $h\getsr \mathcal H$.
%\end{definition}
%
%\begin{definition}[Approximate strongly 2-universal hash function \cite{FOCS:Zuckerman91,C:CacMau97}]
%\label{def:ahash}
%The definition of \emph{approximate strongly 2-universal hash function family} is exactly the same as \cref{def:shash} except that we replace ``uniformly pairwise independent'' by ``approximate pairwise independent''.
%\end{definition}
%
%
%\begin{definition}[2-universal hash function family]
%\label{def:hash}
%	A set $\mathcal H$ of functions with domain $\mathcal X$ and range $\mathcal Y$ is said to be a \emph{2-universal hash function family} if for any distinct $x_1,x_2 \in \mathcal{X}$, 
%	$$\Pr[g(x_1)=g(x_2)]\leq \frac{1}{|\mathcal Y|}.$$
%\end{definition}

%\section{The Proof of \cref{lem:pairwise}}
%\label{sec:pairwiselem}
%In this section we present the proof of \cref{lem:pairwise}.

%\section{Pairwise Independence and Universal Hash Function}
%\label{sec:pairwisedef}
%In this section, we give the backgrounds used to construct our multi-party NIKE in the bounded storage model.
%\heading{Pairwise independence.}
%\label{sec:pairwisedef}
%In this section, we recall the definitions of uniformly and approximately pairwise independence, (approximately) strongly 2-universal hash function family, and 2-universal hash function family. %and recall their instantiations in \cref{fig:con-hash}.~\footnote{The approximate pairwise independence and approximate 2-universal hash function family are defined in \cite{FOCS:Zuckerman91,C:CacMau97} in an implicit way.}
%We refer the reader to~\cref{sec:pairwise} for the definitions of uniformly and approximately pairwise independence, (approximately) strongly 2-universal hash function family, and 2-universal hash function family, and recall their instantiations in \cref{fig:con-hash}.
%Roughly approximately pairwise independent variables are uniformly pairwise independent conditioned on that they are all distinct.

\begin{definition}[Uniform pairwise independence]
\label{def:upairwise}
A sequence of random variables $x_1,\cdots,x_q$ with alphabet $\mathcal Y$ is uniformly pairwise independent if for any $1\leq i<j\leq q$ and any $\hat y_1, \hat y_2 \in \mathcal Y$, we have
	$\Pr[x_i = \hat y_1 \land x_j = \hat y_2] = 1/|\mathcal{Y}|^2$.
\end{definition}


\begin{definition}[Approximate pairwise independence \cite{FOCS:Zuckerman91,C:CacMau97}]
\label{def:apairwise}
A sequence of random variables $x_1,\cdots,x_q$ with alphabet $\mathcal Y$ is approximately pairwise independent if for any $1\leq i<j\leq q$ and any distinct $\hat y_1, \hat y_2 \in \mathcal Y$, we have
$\Pr[x_i = \hat y_1 \land x_j = \hat y_2] = 1/(|\mathcal{Y}|(|\mathcal{Y}|-1))$.
%The definition of \emph{approximate pairwise independence} is exactly the same as \cref{def:upairwise} except that we replace 
%``$\frac{1}{|\mathcal{Y}|^2}$''
%by
%``$\frac{1}{|\mathcal{Y}|(|\mathcal{Y}|-1)}$''.
\end{definition}


%For any $q\in\N$, a sequence of random variables $\vec x_1,\cdots,\vec x_q$ from a range $\mathcal X$ are pairwise independent if for any $1\leq i<j\leq q$ and any $\hy_1, \hy_2 \in \bin^\nn$ with $\hy_1\neq \hy_2$, we have
%	$$\Pr[\vec x_i = \hy_1 \land \vec x_j = \hy_2] = \frac{1}{|\mathcal{X}|^2}.$$

\begin{definition}[Strongly 2-universal hash function family]
\label{def:shash}
	A set $\mathcal H$ of functions with domain $\mathcal X$ and range $\mathcal Y$ is said to be a \emph{strongly 2-universal hash function family} if for any distinct $x_1,x_2 \in \mathcal{X}$, $h(x_1)$ and $h(x_2)$ are uniformly pairwise independent for $h\getsr \mathcal H$.
\end{definition}

\begin{definition}[Approximate strongly 2-universal hash function \cite{FOCS:Zuckerman91,C:CacMau97}]
\label{def:ahash}
The definition of \emph{approximate strongly 2-universal hash function family} is exactly the same as \cref{def:shash} except that we replace ``uniformly pairwise independent'' by ``approximate pairwise independent''.
\end{definition}


\begin{definition}[2-universal hash function family]
\label{def:hash}
	A set $\mathcal H$ of functions with domain $\mathcal X$ and range $\mathcal Y$ is said to be a \emph{2-universal hash function family} if for any distinct $x_1,x_2 \in \mathcal{X}$, we have
	$\Pr[g(x_1)=g(x_2)]\leq 1/|\mathcal Y|$, where $g \getsr \mathcal H$.
\end{definition}


%%We refer the reader to~\cref{sec:pairwise} for the definitions of uniformly and approximately pairwise independence, (approximately) strongly 2-universal hash function family, and 2-universal hash function family, and recall their instantiations in \cref{fig:con-hash}.
%%Roughly approximately pairwise independent variables are uniformly pairwise independent conditioned on that they are all distinct.
%
%\begin{definition}[Uniform pairwise independence]
%\label{def:upairwise}
%A sequence of random variables $x_1,\cdots,x_q$ with alphabet $\mathcal Y$ are uniformly pairwise independent if for any $1\leq i<j\leq q$ and any $\hat y_1, \hat y_2 \in \mathcal X$, we have
%	$\Pr[x_i = \hat y_1 \land x_j = \hat y_2] = 1/|\mathcal{Y}|^2$.
%\end{definition}
%
%
%\begin{definition}[Approximate pairwise independence \cite{FOCS:Zuckerman91,C:CacMau97}]
%\label{def:apairwise}
%A sequence of random variables $x_1,\cdots,x_q$ with alphabet $\mathcal Y$ are approximate pairwise independent if for any $1\leq i<j\leq q$ and any distinct $\hat y_1, \hat y_2 \in \mathcal Y$, we have
%$\Pr[x_i = \hat y_1 \land x_j = \hat y_2] = 1/|\mathcal{Y}|(|\mathcal{Y}|-1)$.
%%The definition of \emph{approximate pairwise independence} is exactly the same as \cref{def:upairwise} except that we replace 
%%``$\frac{1}{|\mathcal{Y}|^2}$''
%%by
%%``$\frac{1}{|\mathcal{Y}|(|\mathcal{Y}|-1)}$''.
%\end{definition}
%
%
%%For any $q\in\N$, a sequence of random variables $\vec x_1,\cdots,\vec x_q$ from a range $\mathcal X$ are pairwise independent if for any $1\leq i<j\leq q$ and any $\hy_1, \hy_2 \in \bin^\nn$ with $\hy_1\neq \hy_2$, we have
%%	$$\Pr[\vec x_i = \hy_1 \land \vec x_j = \hy_2] = \frac{1}{|\mathcal{X}|^2}.$$
%
%\begin{definition}[Strongly 2-universal hash function family]
%\label{def:shash}
%	A set $\mathcal H$ of functions with domain $\mathcal X$ and range $\mathcal Y$ is said to be a \emph{strongly 2-universal hash function family} if for any distinct $x_1,x_2 \in \mathcal{X}$, $h(x_1)$ and $h(x_2)$ are uniformly pairwise independent for $h\getsr \mathcal H$.
%\end{definition}
%
%\begin{definition}[Approximate strongly 2-universal hash function \cite{FOCS:Zuckerman91,C:CacMau97}]
%\label{def:ahash}
%The definition of \emph{approximate strongly 2-universal hash function family} is exactly the same as \cref{def:shash} except that we replace ``uniformly pairwise independent'' by ``approximate pairwise independent''.
%\end{definition}
%
%
%\begin{definition}[2-universal hash function family]
%\label{def:hash}
%	A set $\mathcal H$ of functions with domain $\mathcal X$ and range $\mathcal Y$ is said to be a \emph{2-universal hash function family} if for any distinct $x_1,x_2 \in \mathcal{X}$, we have
%	$\Pr[g(x_1)=g(x_2)]\leq 1/|\mathcal Y|$.
%\end{definition}

\begin{figure}[htpb]
\begin{center}
\framebox{
\small
  \begin{tabular}{l}	
	\begin{minipage}[t]{0.5\textwidth}
	$\mathcal F_n=\{f( x)=a_1\cdot  x+a_0\mid a_0,a_1\in GF(n)\}$\\
	$\mathcal H_n=\{h( x)=a_1\cdot  x+a_0\mid a_0,a_1\in GF(n)\land a_1\neq 0\}$\\
		$\mathcal G_{\lambda,r}=\{g( x)=\msb_r(a\cdot  x)\mid a\in  GF(2^\lambda)\}$
    \end{minipage}
      \end{tabular}}
\end{center} 
\caption{The definitions of $\mathcal F_n$, $\mathcal H_n$, and $\mathcal G_{\lambda,r}$. In this figure, all operations involved in $f$ and $h$ (respectively, $g$) are within the field $GF(n)$ (respectively, $GF(2^\lambda)$). One can see that the outputs of any function in $\mathcal H_n$ with distinct inputs must be all distinct.}
\label{fig:con-hash}
\end{figure}

\begin{lemma}[\cite{C:CacMau97,FOCS:Zuckerman91}]
\label{lem:cm}
For any finite field $GF(n)$ and any $\lambda,r\in\N$, $\mathcal F_n$ (respectively, $\mathcal H_n$) is a strongly 2-universal hash function family (respectively, approximate strongly 2-universal hash function family), and $\mathcal G_{\lambda,r}$ is a 2-universal hash function family.
\end{lemma}
%In the rest of this section, we say that a sequence of variables $(u_i)_{i\in\Z_q}$ is generated by $\mathcal H_n$ if it is generated as $u_i=h(x_i)$ for all $i\in\Z_q$ and some distinct $x_1,\cdots,x_q\in GF(n)$, where $h\getsr\mathcal H_n$.



%Roughly, approximately pairwise independent variables are uniformly pairwise independent conditioned on that they are all distinct.
%
%\subsection{Background of Bounded Storage Model}
%\heading{Pairwise independence.}
%We first recall the definitions of uniformly and approximately pairwise independence, (approximately) strongly 2-universal hash function family, and 2-universal hash function family, and recall their instantiations in \cref{fig:con-hash}.~\footnote{The approximate pairwise independence and approximate 2-universal hash function family are defined in \cite{FOCS:Zuckerman91,C:CacMau97} in an implicit way.}
%%We refer the reader to~\cref{sec:pairwise} for the definitions of uniformly and approximately pairwise independence, (approximately) strongly 2-universal hash function family, and 2-universal hash function family, and recall their instantiations in \cref{fig:con-hash}.
%%Roughly approximately pairwise independent variables are uniformly pairwise independent conditioned on that they are all distinct.
%
%\begin{definition}[Uniform pairwise independence]
%\label{def:upairwise}
%A sequence of random variables $x_1,\cdots,x_q$ with alphabet $\mathcal Y$ are uniformly pairwise independent if for any $1\leq i<j\leq q$ and any $\hat y_1, \hat y_2 \in \mathcal X$, we have
%	$\Pr[x_i = \hat y_1 \land x_j = \hat y_2] = 1/|\mathcal{Y}|^2$.
%\end{definition}
%
%
%\begin{definition}[Approximate pairwise independence \cite{FOCS:Zuckerman91,C:CacMau97}]
%\label{def:apairwise}
%A sequence of random variables $x_1,\cdots,x_q$ with alphabet $\mathcal Y$ are approximate pairwise independent if for any $1\leq i<j\leq q$ and any distinct $\hat y_1, \hat y_2 \in \mathcal Y$, we have
%$\Pr[x_i = \hat y_1 \land x_j = \hat y_2] = 1/|\mathcal{Y}|(|\mathcal{Y}|-1)$.
%%The definition of \emph{approximate pairwise independence} is exactly the same as \cref{def:upairwise} except that we replace 
%%``$\frac{1}{|\mathcal{Y}|^2}$''
%%by
%%``$\frac{1}{|\mathcal{Y}|(|\mathcal{Y}|-1)}$''.
%\end{definition}
%
%
%%For any $q\in\N$, a sequence of random variables $\vec x_1,\cdots,\vec x_q$ from a range $\mathcal X$ are pairwise independent if for any $1\leq i<j\leq q$ and any $\hy_1, \hy_2 \in \bin^\nn$ with $\hy_1\neq \hy_2$, we have
%%	$$\Pr[\vec x_i = \hy_1 \land \vec x_j = \hy_2] = \frac{1}{|\mathcal{X}|^2}.$$
%
%\begin{definition}[Strongly 2-universal hash function family]
%\label{def:shash}
%	A set $\mathcal H$ of functions with domain $\mathcal X$ and range $\mathcal Y$ is said to be a \emph{strongly 2-universal hash function family} if for any distinct $x_1,x_2 \in \mathcal{X}$, $h(x_1)$ and $h(x_2)$ are uniformly pairwise independent for $h\getsr \mathcal H$.
%\end{definition}
%
%\begin{definition}[Approximate strongly 2-universal hash function \cite{FOCS:Zuckerman91,C:CacMau97}]
%\label{def:ahash}
%The definition of \emph{approximate strongly 2-universal hash function family} is exactly the same as \cref{def:shash} except that we replace ``uniformly pairwise independent'' by ``approximate pairwise independent''.
%\end{definition}
%
%
%\begin{definition}[2-universal hash function family]
%\label{def:hash}
%	A set $\mathcal H$ of functions with domain $\mathcal X$ and range $\mathcal Y$ is said to be a \emph{2-universal hash function family} if for any distinct $x_1,x_2 \in \mathcal{X}$, we have
%	$\Pr[g(x_1)=g(x_2)]\leq 1/|\mathcal Y|$.
%\end{definition}

%\subsection{Pairwise Independence}
%We now recall the definitions of (approximate) pairwise independence, (approximately) strongly 2-universal hash function family, and 2-universal hash function family, and provide their instantiations. \footnote{The approximate pairwise independence and approximate 2-universal hash functions are defined in \cite{FOCS:Zuckerman91,C:CacMau97} in an implicit way.}
%
%\begin{definition}[Uniform pairwise independence]
%\label{def:upairwise}
%For any $q\in\N$, a sequence of random variables $x_1,\cdots,x_q$ with alphabet $\mathcal Y$ are uniformly pairwise independent if for any $1\leq i<j\leq q$ and any $\hat y_1, \hat y_2 \in \mathcal X$, we have
%	$$\Pr[x_i = \hat y_1 \land x_j = \hat y_2] = \frac{1}{|\mathcal{Y}|^2}.$$
%\end{definition}
%
%
%\begin{definition}[Approximate pairwise independence \cite{FOCS:Zuckerman91,C:CacMau97}]
%\label{def:apairwise}
%For any $q\in\N$, a sequence of random variables $x_1,\cdots,x_q$ with alphabet $\mathcal Y$ are approximate pairwise independent if for any $1\leq i<j\leq q$ and any distinct $\hat y_1, \hat y_2 \in \mathcal Y$, we have
%$$\Pr[x_i = \hat y_1 \land x_j = \hat y_2] = \frac{1}{|\mathcal{Y}|(|\mathcal{Y}|-1)}.$$
%%The definition of \emph{approximate pairwise independence} is exactly the same as \cref{def:upairwise} except that we replace 
%%``$\frac{1}{|\mathcal{Y}|^2}$''
%%by
%%``$\frac{1}{|\mathcal{Y}|(|\mathcal{Y}|-1)}$''.
%\end{definition}
%One can verify that when a sequence of variables are approximately pairwise independent, they must be all distinct.
%
%%For any $q\in\N$, a sequence of random variables $\vec x_1,\cdots,\vec x_q$ from a range $\mathcal X$ are pairwise independent if for any $1\leq i<j\leq q$ and any $\hy_1, \hy_2 \in \bin^\nn$ with $\hy_1\neq \hy_2$, we have
%%	$$\Pr[\vec x_i = \hy_1 \land \vec x_j = \hy_2] = \frac{1}{|\mathcal{X}|^2}.$$
%
%\begin{definition}[Strongly 2-universal hash function family]
%\label{def:shash}
%	A set $\mathcal H$ of functions with domain $\mathcal X$ and range $\mathcal Y$ is said to be a \emph{strongly 2-universal hash function family} if for any distinct $x_1,x_2 \in \mathcal{X}$, $h(x_1)$ and $h(x_2)$ are uniformly pairwise independent for $h\getsr \mathcal H$.
%\end{definition}
%
%\begin{definition}[Approximate strongly 2-universal hash function \cite{FOCS:Zuckerman91,C:CacMau97}]
%\label{def:ahash}
%The definition of \emph{approximate strongly 2-universal hash function family} is exactly the same as \cref{def:shash} except that we replace ``uniformly pairwise independent'' by ``approximate pairwise independent''.
%\end{definition}
%
%
%\begin{definition}[2-universal hash function family]
%\label{def:hash}
%	A set $\mathcal H$ of functions with domain $\mathcal X$ and range $\mathcal Y$ is said to be a \emph{2-universal hash function family} if for any distinct $x_1,x_2 \in \mathcal{X}$, 
%	$$\Pr[g(x_1)=g(x_2)]\leq \frac{1}{|\mathcal Y|}.$$
%\end{definition}
%
%The following theorem is a direct combination from the Markov’s and Chebyshev’s Inequality.
%
%\begin{theorem}
%\label{thm:cheb}
%For any $n\in\N$, let $( t_i)_{i\in [n]}$ be a sequence of pairwise independent variables with $\expect[ t_i]=\mu_i$ and $\var[ t_i]=\sigma_i^2$. Then for any $a>0$, we have
%$$\Pr[|\sum_{i=1}^n t_i-\sum_{i=1}^n\mu_i|\geq a]\leq \frac{\sum_{i=1}^n\sigma_i^2}{a^2}.$$
%\end{theorem}
%%For any $q\in\N$, a sequence of random variables $\vec x_1,\cdots,\vec x_q$ from a range $\mathcal X$ are pairwise independent if for any $1\leq i<j\leq q$ and any $\hy_1, \hy_2 \in \bin^\nn$ with $\hy_1\neq \hy_2$, we have
%%	$$\Pr[\vec x_i = \hy_1 \land \vec x_j = \hy_2] = \frac{1}{|\mathcal{X}|^2}.$$
%\begin{definition}[Strongly 2-universal hash function family]
%\label{def:shash}
%	A set $\mathcal H$ of functions with domain $\mathcal X$ and range $\mathcal Y$ is said to be a \emph{strongly 2-universal hash function family} if for any distinct $x_1,x_2 \in \mathcal{X}$, $h(x_1)$ and $h(x_2)$ are uniformly pairwise independent for $h\getsr \mathcal H$.
%\end{definition}
%
%\begin{definition}[Approximate strongly 2-universal hash function \cite{FOCS:Zuckerman91,C:CacMau97}]
%\label{def:ahash}
%The definition of \emph{approximate strongly 2-universal hash function family} is exactly the same as \cref{def:shash} except that we replace ``uniformly pairwise independent'' by ``approximate pairwise independent''.
%\end{definition}

%\begin{definition}[2-universal hash function family]
%\label{def:hash}
%	A set $\mathcal H$ of functions with domain $\mathcal X$ and range $\mathcal Y$ is said to be a \emph{2-universal hash function family} if for any distinct $x_1,x_2 \in \mathcal{X}$, we have
%	$\Pr[g(x_1)=g(x_2)]\leq \frac{1}{|\mathcal Y|}$.
%\end{definition}
%>>>>>>> f06d3e8e3fe5b22531e8d7ef4d174804921c0939



%\begin{definition}[Almost pairwise independence]
%\label{def:ahash}
%The definition of \emph{approximate strongly 2-universal hash function family} is exactly the same as \cref{def:pairwise} except that we replace 
%$$``\Pr[h(x_1) = y_1 \land h(x_2) = y_2] = \frac{1}{|\mathcal{Y}|^2}"$$
%by
%$$``\Pr[h(x_1) = y_1 \land h(x_2) = y_2] = \frac{1}{|\mathcal{Y}|(|\mathcal{Y}|-1)}".$$
%\end{definition}

%The definition of \emph{$\kk$-approximate pairwise independence} is exactly the same as \cref{def:pairwise} except that we replace 
%$$``\Pr[x_i = \hat y_1 \land x_j = \hat y_2] = \frac{1}{2^{2\nn}}"$$
%by
%$$``\Pr[x_i = \hat y_1 \land x_j = \hat y_2] = \frac{1}{2^{2\nn}-2^{2\nn-\kk}}".$$
%\end{definition}
%In the rest of this work, when we use the term ``approximate pairwise independence'', we are specifically referring to ``$\nn$-approximate pairwise independence''.
%
%\begin{definition}[$\kk$-approximate pairwise independence]
%\label{def:ahash}
%The definition of \emph{$\kk$-approximate pairwise indedependence} is exactly the same as \cref{def:pairwise} except that we replace 
%$$``\Pr[x_i = y_1 \land x_j = y_2] = \frac{1}{2^{2\nn}}."$$
%by
%$$``\Pr[x_i = y_1 \land x_j = y_2] = \frac{1}{2^{\nn}-2^{2\nn-k}}."$$
%\end{definition}
%In the rest of this work, when we use the term ``approximate pairwise independence'', we are specifically referring to ``$\nn$-approximate pairwise independence''.

%	A set $\mathcal H$ of functions with domain $\mathcal{X}$ and range $\mathcal{Y}$ is said to be an \emph{approximate strongly 2-universal hash function family from $\mathcal{X}$ to $\mathcal{Y}$} if for any $x_1,x_2 \in \mathcal{X}$ such that $x_1 \neq x_2$ and $y_1, y_2 \in \mathcal{Y}$, we have
%	$$\Pr[h(x_1) = y_1 \land h(x_2) = y_2] = \frac{1}{|\mathcal{Y}|(|\mathcal{Y}-1|)}$$
%	for $h\getsr\mathcal H$.

%\begin{definition}[$\kk$-approximate strongly 2-universal hash function family]
%\label{def:khash}
%	A set $\mathcal H_\kk$ of functions with domain $\mathcal X$ and range $\mathcal Y=\bin^\nn$ is said to be a \emph{$\kk$-approximate strongly 2-universal hash function family} if for any $x_1,x_2 \in \mathcal X$ such that $\msb_k(x_1) \neq \msb_k(x_2)$ and $y_1, y_2 \in \mathcal Y$, we have
%	%$$\Pr[\msb_k(h(x_1)) = \msb_k(y_1) \land \msb_k(h(x_2)) = \msb_k(y_2)] = \frac{1}{2^{2k\nn}-2^{2\nn\cdot\kk}}$$
%	$$\Pr[h(x_1) = y_1 \land h(x_2) = y_2] = \frac{1}{2^{2\nn}-2^{2\nn\cdot\kk}}$$
%	for $h\getsr\mathcal H_\kk$.
%\end{definition}

%\begin{figure}[htpb]
%\begin{center}
%\framebox{
%\small
%  \begin{tabular}{l}	
%	\begin{minipage}[t]{0.6\textwidth}
%	{$\mathcal H_s=\{h(\vec x)=\msb_\nn(\vec a_1\cdot \vec x+\vec a_0)\mid \vec a_0,\vec a_1\in GF(2^\qq)\}$}\\
%		{$\mathcal H_u=\{h(\vec u)=\msb_\nn(\vec a_1\cdot \vec x)\mid \vec a_1\in GF(2^\qq)\}$}
%    \end{minipage}
%      \end{tabular}}
%\end{center} 
%\caption{The definitions of $\mathcal H_s$ and $\mathcal H_u$. In this figure, all the operations involved are within $GF(\qq)$.}
%\label{fig:hash}
%\end{figure}
%
%\begin{lemma}[\cite{C:CacMau97}]
%\label{lem:cm}
%$H_u$ (respectively, $\mathcal H_s$) is a 2-universal hash function family (respectively, strongly 2-universal hash function family).
%\end{lemma}





%\begin{figure}[htpb]
%\begin{center}
%\framebox{
%\small
%  \begin{tabular}{l}	
%	\begin{minipage}[t]{0.7\textwidth}
%	$\mathcal H_n=\{(h(\vec x)=(\vec a_1\cdot \msb_k(\vec x)+\vec a_0)\circ \msb_{\qq-\kk}(\vec a_3\cdot \vec x+\vec a_2)\mid \\\vec a_0,\vec a_1\in GF(2^\kk), \vec a_1\neq \vec 0, \vec a_2,\vec a_3\in GF(2^{\qq})\}$
%    \end{minipage}
%      \end{tabular}}
%\end{center} 
%\caption{The definition of $\mathcal H_n$. In this figure, the operations in computing $\vec a_1\cdot \vec x+\vec a_0$ and $\vec a_3\cdot \vec x+\vec a_2$ are within $GF(2^\kk)$ or $GF(\qq)$ respectively.}
%\label{fig:hash}
%\end{figure}

%\begin{lemma}[\cite{C:CacMau97}]
%\label{lem:cm}
%Let $q\in[2^n]$ and $\vec x_1, \vec x_2, \cdots, \vec x_q$ be any $q$ distinct elements in $GF(2^\qq)$. Then the sequence $(\vec h(x_1),\cdots,h(\vec x_q))$ where $h\getsr\mathcal H_u$ is pairwise independent.
%\end{lemma}

%\begin{lemma}[\cite{C:CacMau97}]
%\label{lem:cm2}
%Let $q\in[2^n]$ and $\vec x_1, \vec x_2, \cdots, \vec x_q$ be any $q$ distinct elements in $GF(2^\qq)$. Then the sequence $(\vec h(x_1),\cdots,h(\vec x_q))$ where $h\getsr\mathcal H_u$ is pairwise independent.
%\end{lemma}

%\subsection{Instantiations of Universal Hash Functions}
%We now recall instantiations of the above function families in \cref{fig:con-hash}.
%\begin{figure}[htpb]
%\begin{center}
%\framebox{
%\small
%  \begin{tabular}{l}	
%	\begin{minipage}[t]{0.7\textwidth}
%	$\mathcal F_n=\{f( x)=a_1\cdot  x+a_0\mid a_0,a_1\in GF(n)\}$\\
%	$\mathcal H_n=\{h( x)=a_1\cdot  x+a_0\mid a_0,a_1\in GF(n)\land a_0\neq 0\}$\\
%		$\mathcal G_{n,m}=\{g(u)=\msb_m(a\cdot  x)\mid a\in  GF(2^n)\}$
%    \end{minipage}
%      \end{tabular}}
%\end{center} 
%\caption{The definitions of $\mathcal F_n$, $\mathcal H_n$, and $\mathcal G_{n,m}$. In this figure, all operations involved in $f$ and h (respectively, g) are within the field $GF(n)$ (respectively, $GF(2^n)$). One can see that the outputs of any function in $\mathcal H$ with distinct inputs must be all distinct.}
%\label{fig:con-hash}
%\end{figure}
%\begin{lemma}[\cite{C:CacMau97,FOCS:Zuckerman91}]
%\label{lem:cm}
%For any $m\in\N$, $\mathcal F_m$ (respectively, $\mathcal H_m$) is an strongly 2-universal hash function family (respectively, approximate strongly 2-universal hash function family), and $\mathcal G_m$ is a 2-universal strong hash function family.
%\end{lemma}


%\heading{Key extraction from randomly selected subset.}
%We now recall a main theorem in \cite{C:CacMau97} for privacy amplification. 

%which demonstrates the existence of a way to extract a random key from a subset of a random string with indices being pairwise independent with a small amount of memory. 
%
%\begin{figure}[htpb]
%\begin{center}
%\framebox{
%\small
%  \begin{tabular}{l}	
%	\begin{minipage}[t]{0.8\textwidth}
%	$\mathcal H=\{h( x)=\msb_\nn(a_1\cdot  x+a_0)\mid a_0,a_1\in GF(2^\qq)\}$\\
%	$\mathcal H=\{h( x)=\msb_\nn(a_1\cdot  x+a_0)\mid a_0,a_1\in GF(2^\qq), a_1\neq\vec 0\}$\\
%		{$\mathcal G_r=\{g(\vec u)=\msb_\nn(a_1\cdot  x)\mid a_1\in GF(2^\qq)\}$}
%    \end{minipage}
%      \end{tabular}}
%\end{center} 
%\caption{The definitions of $\mathcal H$, $\mathcal H$, $\mathcal G_r$. In this figure, all the operations involved are within $GF(\qq)$.}
%\label{fig:hash}
%\end{figure}
%
%\begin{lemma}[\cite{C:CacMau97}]
%\label{lem:cm}
%$H_u$ (respectively, $\mathcal H$) is an 2-universal hash function family (respectively, strongly 2-universal hash function family).
%\end{lemma}


%\begin{lemma}
%\label{lem:apairwise}
%Let $\hx_1,\hx_2, \cdots,\hx_q$ a sequence of approximate pairwise independent variables with alphabet $[n]$. Then the sequence $(h(\hx_1),\cdots,h(\hx_n))$ where $h\getsr\mathcal H_\kk$ is $\kk$-approximately pairwise independent.
%\end{lemma}
%\begin{proof}
%%Let $q\in[k]$ and $y,x,x'\in GF(k)$ such that $x\neq x'$.  Then for any $a_0,a_1\in GF(k)$, $a_1\cdot x_i+a_0=a_1\cdot x_j+a_0$ holds if and only if $a_1$ is nonzero.~\footnote{This property is also exploited in \cite{C:CacMau97}.}
%%Hence, we have 
%%$$\Pr[a_0\cdot x+a_1=a_0\cdot x'+a_1|a_0,a_1\getsr GF(k)]\leq \frac{1}{k(k-1)}.$$
%%Moreover, according to \cref{lem:cm}, we have 
%%Let $q\in[k]$, and $y,x_1,\cdots,x_q\in GF(k)$ such that $x_1,\cdots,x_q$ are distinct. 
%For any $1\leq i<j\leq q$ and $a_0,a_1\in GF(k)$, $a_1\cdot \hx_i+a_0=a_1\cdot \hx_j+a_0$ holds if and only if $a_1$ is nonzero.~\footnote{This property is also exploited in \cite{C:CacMau97}.}
%Hence, for any $\hy_1,\hy_2\in GF(k)$ and any $1\leq i<j\leq q$, we have 
%\[\begin{split}\Pr[a_0\cdot \msb_\kk(\hx_i)+a_1=\msb_\kk(\hy_1)\land a_0\cdot\msb_\kk(\hx_j)+a_1=\msb_\kk(\hy_2)\mid\\
%a_0,a_1\getsr GF(k)]\leq \frac{1}{k(k-1)}.\end{split}\]
%Moreover, according to \cref{lem:cm}, we have 
%\[\begin{split}\Pr[\msb_\kk(a_3\cdot \hx_i+a_2)=\lsb_\kk(\hy_1)\land \msb_\kk(a_3\cdot\hx_j+a_2)=\lsb_\kk(\hy_2)\mid\\
% a_2,a_3\getsr GF(2^{\qq})]\leq \frac{1}{2^{2(\qq-\kk)}}.\end{split}\]
% 
% Since the above two events are independent, we have
% $$\Pr[h(\hx_i)=h(\hx_j)\mid h\getsr\mathcal H_\kk]\leq \frac{1}{k(k-1)\cdot 2^{2(\qq-\kk)}}=\frac{1}{2^{2\qq}-2^{2\qq-\kk}},$$
%completing the proof of \cref{lem:apairwise}.
%\end{proof}


%\begin{lemma}
%Let $X_1,\cdots,X_n$ be $\kk$-approximately pairwise independent random variables. Then for any $a>0$:
%$$
%\Pr[|\sum^n_{i=1}X_i-\sum^n_{i=1}\expect(X_i)|\geq a]\leq\frac{\sum_{i=1}^n \var(X_i)}{a^2}.
%$$
%\end{lemma}



%\begin{theorem}[\cite{C:CacMau97}]
%\label{thm:cm}
%Let $\lambda,n\in\N$. Let $\A$ be any adversary  with the memory bound $\amem<n$, $\vec s=(s_i)_{i\in[\lambda]}$ be a sequence of approximately pairwise independent variables generated by $\mathcal H_n$, $\urs$, $\vec z$, and $g$ be random variables such that $\urs \getsr\bin^n$, $\vec z=\A(\urs)$, $g_{\lambda,r}\getsr\mathcal G_r$, and $\key=g(\urs_{\vec s})$. There exists a security event $\event$ such that 
%$$\exists \event:\ \Pr[\event]\geq 1-\epsilon_1-\epsilon_2\ \mbox{and}\ I(\key;g|\vec z,\vec s,\event)\leq \amem.\footnote{The original theorem is presented in a stronger form by asserting that it holds for all possible values of $\vec z$ and $\vec s$, while this weaker version is clearly implied the original one and suffices in our context.}$$
%\end{theorem}


%\begin{lemma}[\cite{C:CacMau97,FOCS:Zuckerman91}]
%%There exist strongly 2-universal hash function families, approximately strongly 2-universal hash function families, and 2-universal hash function families with domain $\bin^m$ and range $\bin^n$ for any $m,n\in\N$. Each of the strongly 2-universal
%\end{lemma}


%\subsection{Strongly 2-Universal Hash Function Family}
%\begin{figure}[htpb]
%\begin{center}
%\framebox{
%\small
%  \begin{tabular}{l}	
%	\begin{minipage}[t]{0.95\textwidth}
%	$\mathcal F_{y,z}=\{f(x)=(a_1\cdot x+a_0\mod y)\cdot z+(b_1\cdot  x+b_0\mod z)\mid a_0,a_1\in \Z_y,b_0,b_1\in \Z_z$\\
%	$\mathcal H_{y,z}=\{h( x)=(a_1\cdot  x+a_0\mod y)\cdot z+(b_1\cdot  x+b_0\mod z)\mid a_0,a_1\in \Z_y,b_0,b_1\in \Z_z, a_1\neq 0\ \mbox{or}\ b_1\neq 0\}$
%    \end{minipage}
%      \end{tabular}}
%\end{center} 
%\caption{The definitions of $\mathcal F_{y,z}$ and $\mathcal H_{y,z}$.
%% In this figure, all operations are within $GF(y\cdot z)$.
% }
%\label{fig:con-hash'}
%\end{figure}
%
%%Due to the fact that 
%%$$(a_1\cdot x+a_0\mod y) \cdot z+ (b_1\cdot  x+b_0\mod z)=(a_1\cdot z+b_1)\cdot x+(a_0\cdot z+b_0)\mod (y\cdot z)$$
%%and $a_1\cdot z+b_1$ and $a_0\cdot z+b_0$ are uniformly distributed
%
%
%
%\begin{lemma}
%\label{lem:cr}
%For any $y,z$ such that $\Z_y$ and $\Z_z$ are fields, $\mathcal F_{y,z}$ (respectively, $\mathcal H_{y,z}$) is a strongly 2-universal hash function family (respectively, approximate strongly 2-universal hash function family) from $\Z_{\min\{y,z\}}$ to $\Z_{y\cdot z}$.
%\end{lemma}
%\begin{proof}
%For any $x, x'\in \Z_y$,
%any $t=t_1\cdot z+t_2$, and any $t'=t_1'\cdot z+t_2'$ where $t_1,t_1'\in\Z_y$ and $t_2,t_2'\in\Z_z$, one can easily check that
%\[\begin{split}
%&\Pr[f(x)= t\land f(x')= t']
%%=&\Pr[(a_1\cdot x+a_0)\cdot z+ (b_1\cdot  x+b_0)=t]\cdot 
%%\Pr[(a_1\cdot x'+a_0)\cdot z+ (b_1\cdot  x'+b_0)=t']\\
%%=&\Pr[(a_1\cdot x+a_0\mod y) \cdot z+ (b_1\cdot  x+b_0\mod z)=t]\cdot \\
%%&\Pr[(a_1\cdot x'+a_0\mod y) \cdot z+ (b_1\cdot  x'+b_0\mod z)=t']\\
%%=&\Pr[a_1\cdot x+a_0\equiv t_1\mod y, b_1\cdot  x+b_0\equiv t_2\mod z]\cdot \\
%%&\Pr[a_1\cdot x'+a_0\equiv t_2\mod y, b_1\cdot  x'+b_0\equiv t_2'\mod z]\\
%%=\frac{1}{y^2}\cdot\frac{1}{z^2} %\tab (\cref{lem:cm})\\
%=\frac{1}{(y\cdot z)^2}
%\end{split}\]
%where $f\getsr\mathcal F_{y,z}$ and $(a_0,a_1,b_0,b_1)$ are the corresponding randomness, i.e., $\mathcal F_{y,z}$ is a strongly 2-universal hash function family.
%
%Also, since $f( x)=f( x')$ if and only if $a_1=b_1=0$, which happens with probability $\frac{1}{y\cdot z}$, 
%conditioned on $ t\neq  t'$ we have
%\[\begin{split}
%&\Pr[f( x)= t\land f( x')= t']\\
%=&\Pr[f( x)= t\land f( x')= t', f\in\mathcal H_{y,z})]\\
%=&\Pr[f( x)= t\land f( x')= t'\mid f\in\mathcal H_{y,z})]\cdot \Pr[f\in\mathcal H_{y,z}]\\
%=&\Pr[f( x)= t\land f( x')= t'\mid f\in\mathcal H_{y,z})]\cdot \frac{y\cdot z-1}{y\cdot z}.
%\end{split}\]
%Moreover, since $\Pr[f( x)= t\land f( x')= t']=\frac{1}{(y\cdot z)^2}$, 
%we have 
%$$\Pr[f( x)= t\land f( x')= t'\mid f\in\mathcal H_{y,z})]=\frac{1}{y\cdot z(y\cdot z-1)},$$
%i.e., $\mathcal H_{y,z}$ is an approximate strongly 2-universal hash function family,
%completing the proof of \cref{lem:cr}.
%\end{proof}
%\heading{Remark.}
%We require that $\Z_y$ and $\Z_z$ be fields only for simplicity.
%One can easily see that the constructions in \cref{fig:con-hash'} can be generalized by only requiring that $\Z_y$ and $\Z_z$ are isomorphic to fields.
%We require that the domain of the functions be $\Z_{y,z}$ to achieve approximate pairwise independence. As a special case, we can also set $y$ as $z=1$.

%\begin{lemma}
%\label{lem:pairwise}
%%Let $\{ h_1(j)\}_{j\in\Z_q},\cdots,\{ h_\pnum(j)\}_{j\in\Z_q}$ be sequences of approximate pairwise independent random variables generated as $t_i^j=h_i(j)$ where $h_i\getsr \mathcal H_n$ for all $i\in[\pnum]$ and all $j\in\Z_q$.
%%Let $\{ h_\pnum(j)\}_{j\in\Z_q}$ be a sequence of $\kk$-approximate pairwise independent random variables with size $\nn$.
%Let $( u_i)_{i\in[q]}$ be the sequence $(h_\pnum(i))_{i\in[q]}$ such that $ u_i= h_\pnum(i)$ if there are indices $i_1,\cdots,i_{\pnum-1}$ satisfying
%$ h_1(i_1)=\cdots= h_{\pnum-1}(i_{\pnum-1})= h_\pnum(i)$ and $ u_i=\bot$ otherwise. 
%Let $\set$ be the set of indices $\Z_q$ restricted to $u_i\neq 0$.
%Then the distribution of $\vec u_\set$ is identical to $(h_\pnum(j))_{j\in\set}$.
%%Then the sequence $(s_i)_{i\in[q]}$ restricted to those positions different from $\bot$ is pairwise independent.
%\end{lemma}
%
%
%\begin{proof}
%%Let $ t_i^j=h_i(j)$ for all $i\in[\pnum]$ and $j\in\Z_q$ and $( u_j)_{j\in\Z_q}$ be the sequence such that $ u_j= h_1(j)$ if there exist indices $j_2,\cdots,j_{\pnum}$ satisfying
%%$\msb_\kk( h_1(j))=\msb_\kk( t_i^{j_2})=\cdots=\msb_\kk( t_\pnum^{j_\pnum})$ and $ u_i=\bot$ otherwise.
%Let $\hx_1,\hx_2$ be any distinct element in $GF(k)$.
%To prove this lemma, it is sufficient to show that 
%%the sequence $( u_i)_{i\in[q]}$ restricted to $ u_i\neq\bot$ is $\kk$-approximately pairwise independent, i.e. $( u_i)_{i\in[q]}$ satisfies
%$
%\Pr[ u_i=\hx_1\land  u_j=\hx_2| u_i\neq \bot\land  u_j\neq\bot]=\frac{1}{k(k-1)}
%$
%for all $i,j\in\Z_q$. 
%
%Due to approximately pairwise independence, we have 
%$
%\Pr[ h_\pnum(i)=\hx_1\land  h_\pnum(j)=\hx_2]=\frac{1}{n(n-1)}
%$
%for all $i,j\in\Z_q$,
%and
%$
%\Pr[\exists i_l,j_l: h_l(i_l)=\hx_1\land h_l(j_l)=\hx_2]=\frac{q(q-1)}{n(n-1)}
%$
%for all $l\in[\pnum-1]$.
%Hence, for any $i,j\in\Z_q$, we have
%\[\begin{split}
%&\Pr[ u_i=\hx_1\land  u_j=\hx_2]\\
%=&\Pr[\exists (i_l,j_l)_{l\in[\pnum-1]}: 
% h_1(i_1)=\cdots= h_{\pnum-1}(i_{\pnum-1})=\hx_1= h_\pnum(i)\land h_1(j_1)=\cdots= h_{\pnum-1}(j_{\pnum-1})=\hx_2= h_\pnum(j)]\\
%=&\Pr[ h_\pnum(i)=\hx_1\land  h_1(j)=x_2]\cdot\prod_{l=1}^\pnum\Pr[\exists i_l,j_l: h_l(i_1)=\hx_1\land  h_l(j_l)=\hx_2]\\
%%=\Pr[&t_1^i=x_1\land h_1(j)=x_2]\cdot\Pr[\exists i_2,\cdots,i_\pnum,j_2,\cdots,j_\pnum: \\&
%%\msb_\kk(t_2^{i_2})=\cdots=\msb_\kk(t_\pnum^{i_\pnum})=\msb_\kk(x_1)\land\\
%%&\msb_\kk(t_2^{j_2})=\cdots=\msb_\kk(t_\pnum^{j_\pnum})=\msb_\kk(x_2)]\\
%=&\frac{1}{k(k-1)}\cdot \Big(\frac{q(q-1)}{n(n-1)}\Big)^{\pnum-1}
%\end{split}\]
%and 
%\[\begin{split}
%&\Pr[ u_i\neq\bot\land  u_j\neq \bot]
%\\
%=&\Pr[\exists (i_l,j_l)_{l\in[\pnum-1]}: h_1(i_1)=\cdots= h_{\pnum-1}(i_{\pnum-1})= h_\pnum(i)\land  h_1(j_1)=\cdots= h_{\pnum-1}(j_{\pnum-1})= h_\pnum(j)]
%\\
%=&\prod_{l=1}^\pnum\Pr[\exists i_l,j_l: h_l(i_l)=h_\pnum(i)\land  h_l(j_l)= h_\pnum(j)]
%=\Big(\frac{q(q-1)}{n(n-1)}\Big)^{\pnum-1}
%\end{split}\]
%i.e.,
%$$
%\Pr[ u_i=x_1\land  u_j=x_2| u_i\neq \bot\land  u_j\neq\bot]=\frac{\Pr[ u_i=\hx_1\land  u_j=\hx_2]}{\Pr[ u_i\neq\bot\land  u_j\neq \bot]}=\frac{1}{k(k-1)}.
%$$
%Moreover, for each other index  $l\in\set$ we have $u_l=h_\pnum(l)$ where $h_\pnum$ can be determined by $u_i$ and $u_j$. Hence, the distribution of $(u_i)_{i\in\set}$ is identical to that of $(h_\pnum(i))_{i\in\set}$, completing the proof of \cref{lem:pairwise}.
%\end{proof}

%\subsection{}
%\subsection{Our Construction}
%\label{sec:bsmke-con}

%Before giving our construction, we recall following theorem, which is a direct combination from the Markov’s and Chebyshev’s Inequality.
%%We recall the following theorem, which is a direct combination of Markov’s and Chebyshev’s Inequalities.
%
%
%\begin{theorem}
%\label{thm:cheb}
%For any $n\in\N$, let $( t_i)_{i\in [n]}$ be a sequence of pairwise independent variables. Then for any $a>0$, we have
%$$\Pr[|\sum_{i=1}^n t_i-\sum_{i=1}^n\expect[ t_i]|\geq a]\leq \frac{\sum_{i=1}^n\var[ t_i]}{a^2}.$$
%\end{theorem}
%
%We also refer the readers to \cref{sec:pairwisedef} for the definitions of uniformly and approximately pairwise independence, (approximately) strongly 2-universal hash function family, and 2-universal hash function family, and their instantiations.

%data processing inequality


%and $\widehat{\vec z}$ and $\widehat{\vec s}$ be any possible values of $\vec z$ and $\vec s$ respectively
%Next we prove the following corollary of the above theorem, regarding to a slightly different distribution of the pairwise independent sequence $\vec s$.
%\begin{corollary}
%\label{cor:cm}
%Let $\lambda,n\in\N$, $\lambda<k\leq n$. Let $\A$ be any adversary with memory bound $\amem<n$, $\urs\getsr\bin^n$, $\vec z\getsr \A(\urs)$, and $\vec s=(s_i)_{i\in[\lambda]}$ be a sequence of random variables such that each $s_i$ is generated as $s_i=h(x_i)\cdot\frac{n}{n}+f(x_i)$ where $h\getsr\mathcal H_n$ and $f\getsr\mathcal F_{n/k}$
% for some sequence $x_1,\cdots,x_\lambda\in\Z_q$, $g\getsr\mathcal G_r$, and $\key=g(\urs_{\vec s})$, and $\widehat{\vec z}$ and $\widehat{\vec s}$ be any possible values of $\vec z$ and $\vec s$ respectively. There exists a security event $\event_1$ such that 
%$$\exists \event_1:\ \Pr[\event_1]\geq 1-\epsilon_1-\epsilon_2-\frac{1}{n}\ \mbox{and}\ I(\key;g|\vec z=\widehat{\vec z},\vec s=\widehat{\vec s},\event_1))\leq \amem.$$
%\end{corollary}

%\begin{proof}
%%Let $\lambda<\min\{\log n,\log n-\log n\}$ and each $s_{i\in[\lambda]}$ be generated as $s_i=h_1(i)\circ h_2(i)$ where $h_1\circ h_2\getsr \mathcal H_{\log n,\log n-\log n}$ in the above theorem. 
%Let $h$ be generated by $h\getsr\mathcal H_{k,n/k}$ instead, where the corresponding random coins are $a_0$, $a_1$, $b_0$, and $b_1$ (see \cref{fig:con-hash'}).
%Let $\event_s$ be the event that $a_1\neq 0$. We have $\Pr[\bar\event_s]\leq \frac{1}{n}$. Thus, combining this with \cref{thm:cm}, we have
%$$\exists \event_0:\ \Pr[\event_0\land \event_s]=\Pr[\event_0]-\Pr[\event_0\land\bar\event_s]\geq\Pr[\event_0]-\Pr[\bar\event_s] \geq 1-\epsilon_1-\epsilon_2-\frac{1}{n}$$
%and
%$$I(\key;g|\vec z=\widehat{\vec z},\vec s=\widehat{\vec s},\event_1)=I(\key;g|\vec z=\widehat{\vec z},\vec s=\widehat{\vec s},(\event_0,\event_s))\leq \amem,$$
%where $\widehat{\vec z}$ and $\widehat{\vec s}$ be any possible values of $\vec z$ and $\vec s$ when $\event_s$ happens. Here notice that $\event_s\land (\vec s=\widehat{\vec s})$ is equivalent to $\vec s=\widehat{\vec s}$.
%Considering $\event_0\land \event_s$ as a single event $\event_1$, \cref{cor:cm} immediately follows.
%\end{proof}
%



%\subsection{Our Construction}
\heading{Parameters.} Next we follow \cite{C:CacMau97} to define several parameters.
Let $\lambda\in\N$ be the security parameter and $\pnum$ be the number of parties.
Let $\amem,n$ be integers such that $\amem<n$. We define $q$, $\epsilon_1$, $\epsilon_2$, $r$, and $\theta$ satisfying the following equations:
\begin{compactitem}
%\item $\gamma=(1-\frac{1}{n})^\pnum$,
\item $\delta=\min\{0.9453,(1/n)(n-\amem-\log(1/\epsilon_1))\}$,
\item $\rho$ such that $H_b(\rho)+\rho\log(1/\delta)+1/n=\delta$,
\item $\lambda=(q/2)\cdot (q/n)^{\pnum-1}=\lfloor(\rho\epsilon_2^2-\rho2^{-\rho n\log(1/\delta)-2})^{-1}\rfloor$,
\item $r=\lfloor\log\theta+\rho\lambda/2-1\rfloor$,
\end{compactitem}
where $H_b$ is the \emph{binary entropy function} defined as $H_b(\rho) = -\rho \log \rho - (1-\rho)\log(1-\rho).$
As noted in \cite{C:CacMau97}, when $n>>\amem$, we have $\delta=0.9453$, $\rho=1/3$,  $\lambda\approx3/\epsilon_2^2$, and $r\approx \lambda/6$.

Before giving our construction, we recall the following two theorems. The first one is the main theorem in \cite{C:CacMau97} for privacy amplification, and the second one is a direct combination of Chebyshev's inequality and Markov's inequality.
%We recall the following theorem, which is a direct combination of Markov’s and Chebyshev’s Inequalities.

%\heading{Privacy amplification.} We now recall . 
\begin{theorem}[\cite{C:CacMau97}]
\label{thm:cm}
Let $\lambda\in\N$ and $GF(n)$ be a finite field. Let $\A$ be any adversary with the memory bound $\amem<n$, $\urs \getsr\bin^n$, $\vec z=\A(\urs)$, $g\getsr\mathcal G_{\lambda,r}$, $\vec s=(s_i)_{i\in[\lambda]}$ be a sequence of approximate pairwise independent
indices over $GF(n)$, and $\key=g(\urs_{\vec s})$. There exists a security event $\event$ such that
$$\exists \event:\ \Pr[\event]\geq 1-\epsilon_1-\epsilon_2\ \mbox{and}\ I(\key;g,\vec z,\vec s\mid\event)\leq \theta.\footnote{The original theorem asserts that $I(\key;g\mid\vec z,\vec s,\event)\leq\theta$, %and the theorem holds for all possible values of $\vec z$ and $\vec s$,
while their proof clearly implies our version.
%this weaker version is clearly implied the original one and suffices in our context.
}$$
\end{theorem}

\begin{theorem}
\label{thm:cheb}
For any $n\in\N$, let $( t_i)_{i\in [n]}$ be a sequence of pairwise independent variables. Then for any $a>0$, we have
$$\Pr[|\sum_{i=1}^n t_i-\sum_{i=1}^n\expect[ t_i]|\geq a]\leq \frac{\sum_{i=1}^n\var[ t_i]}{a^2}.$$
\end{theorem}

%%We now recall a main theorem in \cite{C:CacMau97} for privacy amplification. 
%\heading{Privacy amplification.} We now recall the main theorem in \cite{C:CacMau97} for privacy amplification. 
%\begin{theorem}[\cite{C:CacMau97}]
%\label{thm:cm}
%Let $\lambda,n\in\N$. Let $\A$ be any adversary  with the memory bound $\amem<n$, $\vec s=(s_i)_{i\in[\lambda]}$ be a sequence of approximately pairwise independent variables generated by $\mathcal H_n$, $\urs$, $\vec z$, and $g$ be random variables such that $\urs \getsr\bin^n$, $\vec z=\A(\urs)$, $g_{\lambda,r}\getsr\mathcal G_r$, and $\key=g(\urs_{\vec s})$. There exists a security event $\event$ such that 
%$$\exists \event:\ \Pr[\event]\geq 1-\epsilon_1-\epsilon_2\ \mbox{and}\ I(\key;g|\vec z,\vec s,\event)\leq \theta.\footnote{The original theorem is presented in a stronger form by asserting that it holds for all possible values of $\vec z$ and $\vec s$, while this weaker version is clearly implied the original one and suffices in our context.}$$
%\end{theorem}


%\heading{Construction.} Let $\mathcal H_n$ be the approximately strongly 2-universal hash function family, and $\mathcal G_{\lambda,r}$ be the 2-universal hash function family, as defined in \cref{fig:con-hash}. Let $\mathcal C_1$ be the class of streaming word-RAM families with memory bound $4\log n+\log\lambda+q$, where $q=n\sqrt[\pnum]{\frac{2\lambda}{n}}$.
%We give our construction of NIKE in the BSM as in \cref{fig:con-bsmke}.
%
%
%\begin{figure}[htpb]
%\begin{center}
%\framebox{
%\small
% \begin{tabular}{l}	
%\begin{minipage}[t]{0.95\textwidth}	
%$\underline{\Gen(\urs\in\bin^n)}$:
%\item Set $\pk_i=(h_i,g_i)$ where $h_i\getsr\mathcal H_n$ and $g_i\getsr\mathcal G_{\lambda,r}$
%\item Set $\sk_i=\urs_{\set_i}$ where $\set_i$ is the set of all bits in $\urs$ with indices in $(h_i(j))_{j\in[q]}$
%\item Return $(\pk_i,\sk_i)$\\
%%\item $\party_\pnum$ samples $h_\pnum\getsr\mathcal H_n$, $f\getsr\mathcal H_{n/k}$ and $g\getsr\mathcal G_r$ and stores $f$, $g$, and $\sk_\pnum=\urs_{\set_\pnum}$ where $\set_\pnum=(h_\pnum(j)\cdot\frac{n}{n}+f(j))_{j\in\Z_q}$
%	
%
%	%\item compute $t_i^j=g_i(j)$ for $j=1,\cdots,q$ and broadcast $\pk_i=(t_i^j)$ in a streaming manner (without storing it)
%	$\underline{\share(i,\sk_i,(\pk_j)_{j\in[\pnum]})}$
%	\item Remove bits in $|\sk_i|$ with indices not in $\set_1\cap\cdots \cap\set_\pnum$ by computing $(h_i(j))_{i\in[\pnum],j\in[q]}$ for all $i$ in a streaming manner
%	\item Abort the whole protocol if $|\sk_i|< \lambda$
%%	\item Each $\party_{i\in[\pnum]}$
%%	\begin{compactitem}
%%	\item for $j=1,\cdots,q$, check whether $h_i(j)=h_{i'}(j')$ (or $h_i(j)=\msb_\kk(h(j'))$) for some $j'\in[q]$
%%	\item if the check fails, remove $\urs_{h_i(j)}^*$ from $\sk_i$
%%	\item  store $h$
%%	\end{compactitem}
%%	\item if $|\sk_i|<\lambda$, abort
%%	\item for each $i=2,\cdots,\pnum$ and $j=1,\cdots,q$, $\party_i$ finds $j'\in[q]$ such that $h_1(j)=h_i(j')$, move the $j'$th bit to the $j$th position in $\sk_i$, and discard $h_1$ //make $\sk_i$ identical to $\sk_1$
%	\item Return $\key=g_\pnum(\sk_i)$
%	
%   \end{minipage}
%	\end{tabular}}
%\end{center}
%\caption{The definition of $\bsmke$. Recall that for any set $\set$, $\urs_\set$ denotes the sequence of all bits in $\urs$ with indices (ordered lexicographically) in $\set$.}
%\label{fig:con-bsmke}
%\end{figure}

\heading{Construction.} Let $\mathcal H_n$ be the approximately strongly 2-universal hash function family, and $\mathcal G_{\lambda,r}$ be the 2-universal hash function family, as defined in \cref{fig:con-hash}. Let $\mathcal C_1$ be the class of streaming word-RAM families with memory bounds $4\log n+\lambda+q$, where $q=n\sqrt[\pnum]{2\lambda/n}$.
Throughout this construction, the parameters satisfy $q\leq n$.
Our construction of multi-party NIKE in the BSM is defined as in \cref{fig:con-bsmke}.

\begin{figure}[htpb]
\begin{center}
\framebox{
\small
 \begin{tabular}{l}	
\begin{minipage}[t]{0.9\textwidth}	
$\underline{\Gen(\urs\in\bin^n)}$:
\item Set $\pk_i=(h_i,g_i)$ where $h_i\getsr\mathcal H_n$ and $g_i\getsr\mathcal G_{\lambda,r}$
\item Set $\sk_i=(h_i,\urs_{\set_i})$ where $\set_i$ is the set of indices in $(h_i(j))_{j\in[q]}$
	\item Return $(\pk_i,\sk_i)$\\
%\item $\party_\pnum$ samples $h_\pnum\getsr\mathcal H_n$, $f\getsr\mathcal H_{n/k}$ and $g\getsr\mathcal G_r$ and stores $f$, $g$, and $\sk_\pnum=\urs_{\set_\pnum}$ where $\set_\pnum=(h_\pnum(j)\cdot\frac{n}{n}+f(j))_{j\in\Z_q}$
	

	%\item compute $t_i^j=g_i(j)$ for $j=1,\cdots,q$ and broadcast $\pk_i=(t_i^j)$ in a streaming manner (without storing it)
	$\underline{\share(i,\sk_i,(\pk_j)_{j\in[\pnum]})}$
	\item Parse $\sk_i=(h_i,\urs_{\set_i})$
	\item Remove bits in $\urs_{\set_i}$ with indices not in $\set_1\cap\cdots \cap\set_\pnum$ by computing $(h_i(j))_{i\in[\pnum],j\in[q]}$ in a streaming manner
	\item For each $j\in[q]$, let $u_j=h_\pnum(j)$ if $h_\pnum(j)\in\set_1\cap\cdots\cap\set_\pnum$, and let $u_j=\bot$ otherwise
	\item Let $\vec s$ be obtained from $(u_j)_{j\in[q]}$ by removing all $\bot$'s while preserving the order of $j$
	\item Abort the whole protocol if $|\vec s|< \lambda$
	%\item Each party recovers $\urs_{\vec{s}_\lambda}$ from its stored indexed bits, since $\vec{s}_\lambda\subseteq\set_i$ for every $i\in[\pnum]$
	%	\item Each $\party_{i\in[\pnum]}$
%	\begin{compactitem}
%	\item for $j=1,\cdots,q$, check whether $h_i(j)=h_{i'}(j')$ (or $h_i(j)=\msb_\kk(h(j'))$) for some $j'\in[q]$
%	\item if the check fails, remove $\urs_{h_i(j)}^*$ from $\sk_i$
%	\item  store $h$
%	\end{compactitem}
%	\item if $|\sk_i|<\lambda$, abort
%	\item for each $i=2,\cdots,\pnum$ and $j=1,\cdots,q$, $\party_i$ finds $j'\in[q]$ such that $h_1(j)=h_i(j')$, move the $j'$th bit to the $j$th position in $\sk_i$, and discard $h_1$ //make $\sk_i$ identical to $\sk_1$
	\item Return $\key=g_\pnum(\urs_{\vec{s}_\lambda})$ where $\vec{s}_\lambda$ denotes the first $\lambda$ entries of $\vec s$
	
   \end{minipage}
	\end{tabular}}
\end{center}
\caption{The definition of $\bsmke$. Recall that for any set $\set$, $\urs_\set$ denotes the sequence of all bits in $\urs$ with indices (ordered lexicographically) in $\set$. For any sequence $\vec s$, $\urs_{\vec s}$ follows the order of $\vec s$.}
\label{fig:con-bsmke}
\end{figure}
\begin{theorem}
\label{thm:bsmke}
%For any $\epsilon_1,\epsilon_2>0$, any $0<\epsilon<1$ and any $\amem>0$. 
$\bsmke$ is a $\mathcal C_1$-$\pnum$-NIKE with $2/\lambda$-correctness and $(\amem,\epsilon_1+\epsilon_2+2/\lambda,\theta)$-security in the BSM.
\end{theorem}

%\section{Proof of \cref{thm:bsmke}}
%\label{sec:bsmke-proof}


%We now refer the reader to \cref{sec:pairwisedef} for the definitions of uniformly and approximately pairwise independence, (approximately) strongly 2-universal hash function family, and 2-universal hash function family, and their instantiations.
%In this section we now give the proof of \cref{lem:pairwise}.

%We now give the proof of \cref{thm:bsmke}. %See \cref{sec:bsm-nike} for the instantiations of $\mathcal F_n$, $\mathcal H_n$, and $\mathcal G_{n,m}$, and the parameters.
\begin{proof}
\heading{Complexity.}
First we note that each $\party_{i\in[\pnum]}$ consumes $4\log n+\lambda+q$ bits of memory to store $q$ bits in $\urs$, a 2-universal hash function in $\mathcal G_{\lambda,r}$, and at most two strongly 2-universal hash functions in $\mathcal H_n$ for comparison at any point in the protocol.

\heading{Correctness.}
To prove correctness, it suffices to show that the protocol aborts with probability at most $2/\lambda$.

Let $\set_i=\{h_i(j):j\in[q]\}$ for $i\in[\pnum]$, and define
\[
Y_j=
\begin{cases}
1 & \mbox{if } h_\pnum(j)\in \set_1\cap\cdots\cap\set_{\pnum-1},\\
0 & \mbox{otherwise}
\end{cases}
\ \mbox{and}\
Y=\sum_{j=1}^qY_j .
\]
Since $h_\pnum\getsr\mathcal H_n$ is injective on the $q$ queried points, $Y$ is exactly the number of distinct common indices in $\set_1\cap\cdots\cap\set_\pnum$. Thus the protocol aborts iff $Y<\lambda$.

For each fixed $a\in GF(n)$ and each $i\in[\pnum-1]$, $\Pr[a\in\set_i]=q/n$. Since $\set_1,\ldots,\set_{\pnum-1}$ are independent and $h_\pnum(j)$ is uniform and independent of them,
\[
\expect[Y]
=\sum_{j=1}^q\expect[Y_j]=
\sum_{j=1}^q \Pr[h_\pnum(j)\in \set_1\cap\cdots\cap\set_{\pnum-1}]
=
\sum_{j=1}^q \left(\frac qn\right)^{\pnum-1}
=
q\left(\frac qn\right)^{\pnum-1}
=
2\lambda .
\]

We next bound the variance. For $j\ne j'$,
\[
\begin{aligned}
&\expect[Y_jY_{j'}]=\Pr[Y_j=Y_{j'}=1]\\
=&
\sum_{\substack{a,b\in GF(n)\\a\ne b}}
\Pr[h_\pnum(j)=a,h_\pnum(j')=b]\,
\Pr[a,b\in \set_1\cap\cdots\cap\set_{\pnum-1}] \\
=&
\sum_{\substack{a,b\in GF(n)\\a\ne b}}
\frac{1}{n(n-1)}
\Pr[a,b\in \set_1\cap\cdots\cap\set_{\pnum-1}]
\
(\because \text{approximate pairwise independence of $h_\pnum$})\\
=&
\sum_{\substack{a,b\in GF(n)\\a\ne b}}
\frac{1}{n(n-1)}
\prod_{i=1}^{\pnum-1}\Pr[a,b\in\set_i]
\
(\because \mbox{independence of $(\set_i)_{i=1}^{\pnum-1}$}) \\
=&
\sum_{\substack{a,b\in GF(n)\\a\ne b}}
\frac{1}{n(n-1)}
\left(\frac{q(q-1)}{n(n-1)}\right)^{\pnum-1} \ (\because \mbox{injectivity and uniformity of $\{h_i\}_{i=1}^{\pnum-1}$})\\
=&
\left(\frac{q(q-1)}{n(n-1)}\right)^{\pnum-1}\le
\left(\frac qn\right)^{2(\pnum-1)}
=
\expect[Y_j]\expect[Y_{j'}].
\end{aligned}
\]
%where in the third equality we use that, for every fixed distinct $a,b\in GF(n)$ and every
%$i\in[\pnum-1]$, injectivity and uniformity of $h_i$ imply
%$\Pr[a,b\in\set_i]=q(q-1)/(n(n-1))$.
As a result, we have
\[
\begin{aligned}
\var[Y]
&=
\sum_{j=1}^q\var[Y_j]
+
2\sum_{1\leq j<j'\leq q}
\left(\expect[Y_jY_{j'}]-\expect[Y_j]\expect[Y_{j'}]\right) 
\leq
\sum_{j=1}^q\var[Y_j]
\leq
\sum_{j=1}^q\expect[Y_j]
=
2\lambda.
\end{aligned}
\]
By Chebyshev's inequality,
\[
\Pr[Y<\lambda]
\leq
\Pr[|Y-\expect[Y]|\geq\lambda]
\leq
\frac{\var[Y]}{\lambda^2}
\leq
\frac{2}{\lambda}.
\]
Therefore, the protocol aborts with probability at most $2/\lambda$.
%\heading{Correctness.}
%To prove correctness, it suffices to show that the protocol aborts with probability at most $2/\lambda$.
%
%Let $\set_i=\{h_i(j):j\in[q]\}$ for $i\in[\pnum]$, let
%\[
%Y_j=
%\begin{cases}
%1 & \mbox{if } h_\pnum(j)\in \set_1\cap\cdots\cap\set_{\pnum-1},\\
%0 & \mbox{otherwise},
%\end{cases}
%\]
%and let $Y=\sum_{j=1}^qY_j$. Since $h_\pnum\getsr\mathcal H_n$ is injective on the $q$ queried points, $Y$ is exactly the number of distinct common indices in $\set_1\cap\cdots\cap\set_\pnum$. Thus the protocol aborts iff $Y<\lambda$.
%
%For each fixed $a\in GF(n)$ and each $i\in[\pnum-1]$, we have $\Pr[a\in\set_i]=q/n$. Since $\set_1,\ldots,\set_{\pnum-1}$ are independent and $h_\pnum(j)$ is uniform and independent of them, we have $\expect[Y_j]=(q/n)^{\pnum-1}$. Hence, we have
%$$\expect[Y]=\sum_{j=1}^q\expect[Y_j]=q(q/n)^{\pnum-1}=2\lambda.$$
%
%We now bound the variance. For $j\ne j'$, the pair $(h_\pnum(j),h_\pnum(j'))$ is uniform over ordered pairs of distinct field elements and independent of $\set_1,\ldots,\set_{\pnum-1}$. Hence, we have
%\[
%\Pr[Y_j=Y_{j'}=1]
%=
%\sum_{\substack{a,b\in GF(n),a\ne b}}
%\Pr[h_\pnum(j)=a,h_\pnum(j')=b]\,
%\Pr[a,b\in \set_1\cap\cdots\cap\set_{\pnum-1}].
%\]
%For fixed distinct $a,b\in GF(n)$ and $i\in[\pnum-1]$, injectivity and uniformity of $h_i$ give
%$$\Pr[a,b\in\set_i]=q(q-1)/(n(n-1)).$$ Then, by independence, we have
%\[
%\Pr[Y_j=Y_{j'}=1]
%=
%\left(q(q-1)/(n(n-1))\right)^{\pnum-1}.
%\]
%Since $q\leq n$, we have $q(q-1)/(n(n-1))\leq(q/n)^2$, and therefore
%\[
%\expect[Y_jY_{j'}]
%=
%\Pr[Y_j=Y_{j'}=1]
%\leq
%(q/n)^{2(\pnum-1)}
%=
%\expect[Y_j]\expect[Y_{j'}].
%\]
%Thus the covariance terms are non-positive. Expanding the variance,
%\[
%\begin{aligned}
%\var[Y]
%&=
%\sum_{j=1}^q\var[Y_j]
%+
%2\sum_{1\leq j<j'\leq q}
%\left(\expect[Y_jY_{j'}]-\expect[Y_j]\expect[Y_{j'}]\right) \\
%&\leq
%\sum_{j=1}^q\var[Y_j]
%\leq
%\sum_{j=1}^q\expect[Y_j]
%=
%2\lambda.
%\end{aligned}
%\]
%By Chebyshev's inequality,
%\[
%\Pr[Y<\lambda]
%\leq
%\Pr[|Y-\expect[Y]|\geq\lambda]
%\leq
%\var[Y]/\lambda^2
%\leq
%2/\lambda.
%\]
%Therefore, the protocol aborts with probability at most $2/\lambda$.
%\heading{Correctness.} To prove correctness, it is sufficient to prove that the probability that the protocol aborts is at most $2/\lambda$.
%
%Let $\set_i=\{h_i(j):j\in[q]\}$ for $i\in[\pnum]$, $Y_j=\begin{cases}
%1 \ \mbox{if}\ h_\pnum(j)\in \set_1\cap\cdots\cap\set_{\pnum-1} \\0\ \mbox{otherwise}
%\end{cases}$
%for all $j\in[q]$,
%and $Y=\sum_{j=1}^qY_j$.
%Since $h_\pnum\getsr\mathcal H_n$ is injective on the $q$ queried points, the values $h_\pnum(1),\ldots,h_\pnum(q)$ are distinct. Hence, $Y$ is exactly the number of distinct common indices in $\set_1\cap\cdots\cap\set_\pnum$.
%
%For each fixed $a\in GF(n)$ and every $i\in[\pnum-1]$, we have
%$
%\Pr[a\in\set_i]=q/n.
%$
%Moreover, the sets $\set_1,\ldots,\set_{\pnum-1}$ are generated independently, and $h_\pnum(j)$ is uniform over $GF(n)$ and independent of them. Thus, for every $j\in[q]$, we have
%$
%\expect[Y_j]=(q/n)^{\pnum-1}
%$,
%i.e.,
%$
%\expect[Y]=q(q/n)^{\pnum-1}=2\lambda.
%$
%
%For all $j\ne j'$, $(h_\pnum(j),h_\pnum(j'))$ is uniformly distributed over ordered pairs of distinct field elements and is independent of $\set_1,\ldots,\set_{\pnum-1}$. Hence, we have
%\[
%\Pr[Y_j=Y_{j'}=1]
%=
%\sum_{\substack{a,b\in GF(n)\\a\ne b}}
%\Pr[h_\pnum(j)=a,h_\pnum(j')=b]\,
%\Pr[a,b\in \set_1\cap\cdots\cap\set_{\pnum-1}].
%\]
%For all fixed distinct $a,b\in GF(n)$ and all $i\in[\pnum-1]$, since $h_i\getsr\mathcal H_n$ is injective and uniform on ordered pairs of distinct outputs, we have
%\[
%\Pr[a,b\in\set_i]=q(q-1)/(n(n-1)).
%\]
%Moreover, since the sets $\set_1,\ldots,\set_{\pnum-1}$ are generated independently, we have
%\[
%\Pr[a,b\in \set_1\cap\cdots\cap\set_{\pnum-1}]
%=
%\left(q(q-1)/(n(n-1))\right)^{\pnum-1}.
%\]
%Therefore,
%\[
%\Pr[Y_j=Y_{j'}=1]
%=
%\left(q(q-1)/(n(n-1))\right)^{\pnum-1}.
%\]
%Since $q\leq n$,
%\[
%q(q-1)/(n(n-1))
%\leq
%(q/n)^2,
%\]
%and hence
%\[
%\Pr[Y_j=Y_{j'}=1]
%\leq
%(q/n)^{2(\pnum-1)}
%=
%\expect[Y_j]\expect[Y_{j'}].
%\]
%Since $Y_j$ and $Y_{j'}$ are indicators, $\expect[Y_jY_{j'}]=\Pr[Y_j=Y_{j'}=1]$. Hence the last inequality is equivalent to
%\[
%\expect[Y_jY_{j'}]\leq \expect[Y_j]\expect[Y_{j'}],
%\]
%which is the sense in which the $Y_j$'s are pairwise negatively correlated. Therefore, by expanding the variance directly, we have
%\[
%\var[Y]
%=
%\sum_{j=1}^q\var[Y_j]
%+
%2\sum_{1\leq j<j'\leq q}
%\left(\expect[Y_jY_{j'}]-\expect[Y_j]\expect[Y_{j'}]\right)
%\leq
%\sum_{j=1}^q\var[Y_j]
%\leq
%\sum_{j=1}^q\expect[Y_j]
%=
%\expect[Y]
%=
%2\lambda.
%\]
%Then by Chebyshev's inequality,
%\[
%\Pr[Y<\lambda]
%\leq
%\Pr[|Y-\expect[Y]|\geq \lambda]
%\leq
%\var[Y]/\lambda^2
%\leq
%2/\lambda.
%\]
%Therefore, the protocol aborts with probability at most $2/\lambda$.

%\heading{$(m,\epsilon_1+\epsilon_2),\amem)$-security.}
\heading{Security.}
Let $f:\bin^n\rightarrow\bin^\amem$ be any storage function applied by the adversary. Let $(h_i,g_i)_{i\in[\pnum]}$ be the honestly generated public keys.
 %Let $\{ h_1(j)\}_{j\in\Z_q},\cdots,\{ h_\pnum(j)\}_{j\in\Z_q}$ be sequences of approximate pairwise independent random variables generated as $t_i^j=h_i(j)$ where $h_i\getsr \mathcal H_n$ for all $i\in[\pnum]$ and all $j\in\Z_q$.
%Let $\{ h_\pnum(j)\}_{j\in\Z_q}$ be a sequence of $\kk$-approximate pairwise independent random variables with size $\nn$. 
Let $\vec u=(u_i)_{i\in[q]}$ be the sequence such that $u_i= h_\pnum(i)$ if there are indices $i_1,\cdots,i_{\pnum-1}$ satisfying
$ h_1(i_1)=\cdots= h_{\pnum-1}(i_{\pnum-1})= h_\pnum(i)$ and $ u_i=\bot$ otherwise, and let $\mathcal S$ be the set of all indices of $u_i$ restricted to $u_i\neq \bot$, i.e., $\vec s=\vec u_{\mathcal S}$ be the sequence $\vec u$ restricted to $u_i\neq \bot$.
We have the following lemma.

\begin{lemma}
\label{lem:pairwise}
$\vec s$ is approximately pairwise independent. %\footnote{Recall that as mentioned in \cref{sec:bsmke-con}, we say that $\vec s$ is generated by $\mathcal H_n$ if it is generated as $s_i=h(x_i)$ for all $i\in[n]$ and some distinct $x_1,\cdots,x_n\in GF(2^n)$, where $h\getsr\mathcal H_n$.}
%Then the sequence $(s_i)_{i\in[q]}$ restricted to those positions different from $\bot$ is pairwise independent.
\end{lemma}

\begin{proof}
%Let $ t_i^j=h_i(j)$ for all $i\in[\pnum]$ and $j\in\Z_q$ and $( u_j)_{j\in\Z_q}$ be the sequence such that $ u_j= h_1(j)$ if there exist indices $j_2,\cdots,j_{\pnum}$ satisfying
%$\msb_\kk( h_1(j))=\msb_\kk( t_i^{j_2})=\cdots=\msb_\kk( t_\pnum^{j_\pnum})$ and $ u_i=\bot$ otherwise.
Let $\hx_1,\hx_2$ be any distinct elements in $GF(n)$.
%To prove this lemma, it is sufficient to show that 
%the sequence $( u_i)_{i\in[q]}$ restricted to $ u_i\neq\bot$ is $\kk$-approximately pairwise independent, i.e. $( u_i)_{i\in[q]}$ satisfies
%$
%\Pr[ u_i=\hx_1\land  u_j=\hx_2| u_i\neq \bot\land  u_j\neq\bot]=\frac{1}{k(k-1)}
%$
%for all $i,j\in\Z_q$. 
Due to approximately pairwise independence, we have 
$
\Pr[ h_\pnum(i)=\hx_1\land  h_\pnum(j)=\hx_2]=\frac{1}{n(n-1)}
$
for all $i,j\in[q]$,
and
$
\Pr[\exists i_l,j_l: h_l(i_l)=\hx_1\land h_l(j_l)=\hx_2]=\frac{q(q-1)}{n(n-1)}
$
for all $l\in[\pnum-1]$.
Hence, for any $i,j\in\Z_q$, we have
\[\begin{split}
\Pr[ u_i=\hx_1\land  u_j=\hx_2]
=&\Pr[\exists (i_l,j_l)_{l\in[\pnum-1]}: 
 h_1(i_1)=\cdots= h_{\pnum-1}(i_{\pnum-1})=\hx_1= h_\pnum(i)\land\\
 &\tab h_1(j_1)=\cdots= h_{\pnum-1}(j_{\pnum-1})=\hx_2= h_\pnum(j)]\\
=&\Pr[ h_\pnum(i)=\hx_1\land  h_\pnum(j)=\hx_2]\cdot\prod_{l=1}^{\pnum-1}\Pr[\exists i_l,j_l: h_l(i_l)=\hx_1\land  h_l(j_l)=\hx_2]\\
%=\Pr[&t_1^i=x_1\land h_1(j)=x_2]\cdot\Pr[\exists i_2,\cdots,i_\pnum,j_2,\cdots,j_\pnum: \\&
%\msb_\kk(t_2^{i_2})=\cdots=\msb_\kk(t_\pnum^{i_\pnum})=\msb_\kk(x_1)\land\\
%&\msb_\kk(t_2^{j_2})=\cdots=\msb_\kk(t_\pnum^{j_\pnum})=\msb_\kk(x_2)]\\
=&\frac{1}{n(n-1)}\cdot \Big(\frac{q(q-1)}{n(n-1)}\Big)^{\pnum-1},
\end{split}\]
and
\[\begin{split}
\Pr[ u_i\neq\bot\land  u_j\neq \bot]
=&\Pr[\exists (i_l,j_l)_{l\in[\pnum-1]}: h_1(i_1)=\cdots= h_{\pnum-1}(i_{\pnum-1})= h_\pnum(i)\land\\
&\tab  h_1(j_1)=\cdots= h_{\pnum-1}(j_{\pnum-1})= h_\pnum(j)]
\\
=&\prod_{l=1}^\pnum\Pr[\exists i_l,j_l: h_l(i_l)=h_\pnum(i)\land  h_l(j_l)= h_\pnum(j)]
=\Big(\frac{q(q-1)}{n(n-1)}\Big)^{\pnum-1},
\end{split}\]
i.e.,
$$
 \Pr[ u_i=x_1\land  u_j=x_2| u_i\neq \bot\land  u_j\neq\bot]=\frac{\Pr[ u_i=\hx_1\land  u_j=\hx_2]}{\Pr[ u_i\neq\bot\land  u_j\neq \bot]}=\frac{1}{n(n-1)}.
$$
This completes the proof of \cref{lem:pairwise}.
%Hence, the distribution of $(u_i,u_j)$ is identical to that of $(h_\pnum(i),h_\pnum(j))$, conditioned on neither of them being $\bot$.
%Moreover, for each $l\in\set\backslash\{i,j\}$ we have $u_l=h_\pnum(l)$ where $h_\pnum$ can be determined by $u_i$ and $u_j$. Hence, the distribution of $(u_i)_{i\in\set}$ is identical to that of $(h_\pnum(i))_{i\in\set}$
\end{proof}

By the correctness analysis above, we have
$\Pr[\mathsf{Good}]\geq 1-2/\lambda$. On $\mathsf{Good}$, let
$\vec{s}_{\lambda}$ denote the first $\lambda$ entries of $\vec s$. By
\cref{lem:pairwise} and \cref{thm:cm}, there exists an
event $\event_0$ such that
\[
\Pr[\event_0\mid\mathsf{Good}]\geq 1-\epsilon_1-\epsilon_2
\
\mbox{and}
\
I(\key;g_\pnum,\vec z,\vec{s}_{\lambda}\mid\event_0,\mathsf{Good})
\leq \theta,
\]
where $\urs\getsr\bin^n$, $\vec z=f(\urs)$, and
$\key=g_\pnum(\urs_{\vec{s}_{\lambda}})$.
Let $\event=\event_0\cap\mathsf{Good}$. Then we have
\[
\Pr[\event]\geq 1-\epsilon_1-\epsilon_2-2/\lambda.
\]
The event $\event_0$ is chosen using only $(\urs,g_\pnum,\vec z,\vec{s}_\lambda)$ and not the remaining public-key components. Since $\vec{s}_{\lambda}$ is determined by $(h_i)_{i\in[\pnum]}$ and $\key=g_\pnum(\urs_{\vec{s}_{\lambda}})$, those remaining components give no additional information about $\key$ once $(g_\pnum,\vec z,\vec{s}_{\lambda},\event)$ is fixed. Hence
\[
I(\key;(\vec z,(\pk_i)_{i\in[\pnum]})\mid\event)
=
I(\key;g_\pnum,\vec z,\vec{s}_{\lambda}\mid\event)
\leq
\theta.
\]
Thus, $\bsmke$ satisfies $(\amem,\epsilon_1+\epsilon_2+2/\lambda,\theta)$-security in the BSM.

Putting all the above together, \cref{thm:bsmke} follows.
%\heading{Security.}
%Let $f:\bin^n\rightarrow\bin^\amem$ be any storage function applied by the adversary. Let $(h_i,g_i)_{i\in[\pnum]}$ be the honestly generated public keys.
% %Let $\{ h_1(j)\}_{j\in\Z_q},\cdots,\{ h_\pnum(j)\}_{j\in\Z_q}$ be sequences of approximate pairwise independent random variables generated as $t_i^j=h_i(j)$ where $h_i\getsr \mathcal H_n$ for all $i\in[\pnum]$ and all $j\in\Z_q$.
%%Let $\{ h_\pnum(j)\}_{j\in\Z_q}$ be a sequence of $\kk$-approximate pairwise independent random variables with size $\nn$.
%Let $\vec u=(u_i)_{i\in[q]}$ be the sequence such that $u_i= h_\pnum(i)$ if there are indices $i_1,\cdots,i_{\pnum-1}$ satisfying
%$ h_1(i_1)=\cdots= h_{\pnum-1}(i_{\pnum-1})= h_\pnum(i)$ and $ u_i=\bot$ otherwise, and let $\mathcal S$ be the set of indices $i$ such that $u_i\neq \bot$, i.e., $\vec s=\vec u_{\mathcal S}$ be the sequence $\vec u$ restricted to $u_i\neq \bot$.
%We have the following lemma.
%%We refer the reader to \cref{sec:pairwise} for the proof of \cref{lem:pairwise}.
%\begin{lemma}
%\label{lem:pairwise}
%$\vec s$ is generated by $\mathcal H_n$. %\footnote{Recall that as mentioned in \cref{sec:bsmke-con}, we say that $\vec s$ is generated by $\mathcal H_n$ if it is generated as $s_i=h(x_i)$ for all $i\in[n]$ and some distinct $x_1,\cdots,x_n\in GF(2^n)$, where $h\getsr\mathcal H_n$.}
%%Then the sequence $(s_i)_{i\in[q]}$ restricted to those positions different from $\bot$ is pairwise independent.
%\end{lemma}
%%\section{Multi-party NIKE in the Bounded Storage Model with Background and (Full)Security}
%In this section, we give the definitions and theorems used to construct our multi-party NIKE in the bounded storage model, along with the proof of its security.
%\subsection{Background of Bounded Storage Model}
%\section{Proof of \cref{lem:pairwise}}
%\label{sec:pairwise}

\begin{proof}
%Let $ t_i^j=h_i(j)$ for all $i\in[\pnum]$ and $j\in\Z_q$ and $( u_j)_{j\in\Z_q}$ be the sequence such that $ u_j= h_1(j)$ if there exist indices $j_2,\cdots,j_{\pnum}$ satisfying
%$\msb_\kk( h_1(j))=\msb_\kk( t_i^{j_2})=\cdots=\msb_\kk( t_\pnum^{j_\pnum})$ and $ u_i=\bot$ otherwise.
Let $\hx_1,\hx_2$ be any distinct elements in $GF(n)$.
%To prove this lemma, it is sufficient to show that 
%the sequence $( u_i)_{i\in[q]}$ restricted to $ u_i\neq\bot$ is $\kk$-approximately pairwise independent, i.e. $( u_i)_{i\in[q]}$ satisfies
%$
%\Pr[ u_i=\hx_1\land  u_j=\hx_2| u_i\neq \bot\land  u_j\neq\bot]=\frac{1}{k(k-1)}
%$
%for all $i,j\in\Z_q$. 
Due to approximately pairwise independence, we have 
$$
\Pr[ h_\pnum(i)=\hx_1\land  h_\pnum(j)=\hx_2]=1/(n(n-1))
$$
for all distinct $i,j\in[q]$,
and
$$
\Pr[\exists i_l,j_l: h_l(i_l)=\hx_1\land h_l(j_l)=\hx_2]=q(q-1)/(n(n-1))
$$
for all $l\in[\pnum-1]$.
Hence, for any distinct $i,j\in\Z_q$, we have
\[\begin{split}
\Pr[ u_i=\hx_1\land  u_j=\hx_2]
=&\Pr[\exists (i_l,j_l)_{l\in[\pnum-1]}: 
 h_1(i_1)=\cdots= h_{\pnum-1}(i_{\pnum-1})=\hx_1= h_\pnum(i)\land\\
 &\tab h_1(j_1)=\cdots= h_{\pnum-1}(j_{\pnum-1})=\hx_2= h_\pnum(j)]\\
=&\Pr[ h_\pnum(i)=\hx_1\land  h_\pnum(j)=\hx_2]\cdot\prod_{l=1}^{\pnum-1}\Pr[\exists i_l,j_l: h_l(i_l)=\hx_1\land  h_l(j_l)=\hx_2]\\
%=\Pr[&t_1^i=x_1\land h_1(j)=x_2]\cdot\Pr[\exists i_2,\cdots,i_\pnum,j_2,\cdots,j_\pnum: \\&
%\msb_\kk(t_2^{i_2})=\cdots=\msb_\kk(t_\pnum^{i_\pnum})=\msb_\kk(x_1)\land\\
%&\msb_\kk(t_2^{j_2})=\cdots=\msb_\kk(t_\pnum^{j_\pnum})=\msb_\kk(x_2)]\\
=&1/(n(n-1))\cdot \Big(q(q-1)/(n(n-1))\Big)^{\pnum-1},
\end{split}\]
and
\[\begin{split}
\Pr[ u_i\neq\bot\land  u_j\neq \bot]
=&\Pr[\exists (i_l,j_l)_{l\in[\pnum-1]}: h_1(i_1)=\cdots= h_{\pnum-1}(i_{\pnum-1})= h_\pnum(i)\land\\
&\tab  h_1(j_1)=\cdots= h_{\pnum-1}(j_{\pnum-1})= h_\pnum(j)]
\\
=&\prod_{l=1}^\pnum\Pr[\exists i_l,j_l: h_l(i_l)=h_\pnum(i)\land  h_l(j_l)= h_\pnum(j)]
=\Big(q(q-1)/(n(n-1))\Big)^{\pnum-1},
\end{split}\]
i.e.,
$$
 \Pr[ u_i=\hx_1\land  u_j=\hx_2\mid u_i\neq \bot\land  u_j\neq\bot]=\frac{\Pr[ u_i=\hx_1\land  u_j=\hx_2]}{\Pr[ u_i\neq\bot\land  u_j\neq \bot]}=1/(n(n-1)).
$$
Therefore, conditioned on two surviving positions, their values are uniformly distributed over ordered pairs of distinct field elements. Restricting to the non-$\bot$ entries and then taking the first $\lambda$ entries preserves this pairwise distribution. Hence, conditioned on $\mathsf{Good}$, $\vec{s}_{\lambda}$ satisfies the condition required by \cref{thm:cm}.
\end{proof}


%In this section, we give the backgrounds used to construct our multi-party NIKE in the bounded storage model.
%\heading{Pairwise independence.}
%\label{sec:pairwisedef}
%In this section, we recall the definitions of uniformly and approximately pairwise independence, (approximately) strongly 2-universal hash function family, and 2-universal hash function family, and recall their instantiations in \cref{fig:con-hash}.~\footnote{The approximate pairwise independence and approximate 2-universal hash function family are defined in \cite{FOCS:Zuckerman91,C:CacMau97} in an implicit way.}
%%We refer the reader to~\cref{sec:pairwise} for the definitions of uniformly and approximately pairwise independence, (approximately) strongly 2-universal hash function family, and 2-universal hash function family, and recall their instantiations in \cref{fig:con-hash}.
%%Roughly approximately pairwise independent variables are uniformly pairwise independent conditioned on that they are all distinct.
%
%\begin{definition}[Uniform pairwise independence]
%\label{def:upairwise}
%A sequence of random variables $x_1,\cdots,x_q$ with alphabet $\mathcal Y$ are uniformly pairwise independent if for any $1\leq i<j\leq q$ and any $\hat y_1, \hat y_2 \in \mathcal X$, we have
%	$\Pr[x_i = \hat y_1 \land x_j = \hat y_2] = 1/|\mathcal{Y}|^2$.
%\end{definition}
%
%
%\begin{definition}[Approximate pairwise independence \cite{FOCS:Zuckerman91,C:CacMau97}]
%\label{def:apairwise}
%A sequence of random variables $x_1,\cdots,x_q$ with alphabet $\mathcal Y$ are approximate pairwise independent if for any $1\leq i<j\leq q$ and any distinct $\hat y_1, \hat y_2 \in \mathcal Y$, we have
%$\Pr[x_i = \hat y_1 \land x_j = \hat y_2] = 1/|\mathcal{Y}|(|\mathcal{Y}|-1)$.
%%The definition of \emph{approximate pairwise independence} is exactly the same as \cref{def:upairwise} except that we replace 
%%``$\frac{1}{|\mathcal{Y}|^2}$''
%%by
%%``$\frac{1}{|\mathcal{Y}|(|\mathcal{Y}|-1)}$''.
%\end{definition}
%
%
%%For any $q\in\N$, a sequence of random variables $\vec x_1,\cdots,\vec x_q$ from a range $\mathcal X$ are pairwise independent if for any $1\leq i<j\leq q$ and any $\hy_1, \hy_2 \in \bin^\nn$ with $\hy_1\neq \hy_2$, we have
%%	$$\Pr[\vec x_i = \hy_1 \land \vec x_j = \hy_2] = \frac{1}{|\mathcal{X}|^2}.$$
%
%\begin{definition}[Strongly 2-universal hash function family]
%\label{def:shash}
%	A set $\mathcal H$ of functions with domain $\mathcal X$ and range $\mathcal Y$ is said to be a \emph{strongly 2-universal hash function family} if for any distinct $x_1,x_2 \in \mathcal{X}$, $h(x_1)$ and $h(x_2)$ are uniformly pairwise independent for $h\getsr \mathcal H$.
%\end{definition}
%
%\begin{definition}[Approximate strongly 2-universal hash function \cite{FOCS:Zuckerman91,C:CacMau97}]
%\label{def:ahash}
%The definition of \emph{approximate strongly 2-universal hash function family} is exactly the same as \cref{def:shash} except that we replace ``uniformly pairwise independent'' by ``approximate pairwise independent''.
%\end{definition}
%
%
%\begin{definition}[2-universal hash function family]
%\label{def:hash}
%	A set $\mathcal H$ of functions with domain $\mathcal X$ and range $\mathcal Y$ is said to be a \emph{2-universal hash function family} if for any distinct $x_1,x_2 \in \mathcal{X}$, we have
%	$\Pr[g(x_1)=g(x_2)]\leq 1/|\mathcal Y|$.
%\end{definition}

%We now recall a main theorem in \cite{C:CacMau97} for privacy amplification. 

%%\subsection{The Proof of \cref{lem:pairwise}}
%%\section{The Proof of \cref{lem:pairwise}}
%%\label{sec:pairwise}
%
%Let $\mathsf{Good}$ denote the non-abort event, i.e., $|\vec s|\geq \lambda$.
%By the correctness analysis above, $\Pr[\mathsf{Good}]\geq 1-2/\lambda$.
%Conditioned on $\mathsf{Good}$, $\vec{s}_{\lambda}$ is well-defined.
%Combining \cref{lem:pairwise} with \cref{thm:cm}, there exists an event $\event_0$ such that
%$\Pr[\event_0\mid\mathsf{Good}]\geq 1-\epsilon_1-\epsilon_2$ and
%\[
%I(\key;g_\pnum,\vec z,\vec{s}_{\lambda}\mid\event_0,\mathsf{Good})
%\leq \theta,
%\]
%where $\urs \getsr\bin^n$, $\vec z=f(\urs)$, $\vec{s}_{\lambda}$ denotes the first $\lambda$ entries of $\vec s$, and $\key=g_\pnum(\urs_{\vec{s}_{\lambda}})$.
%Let $\event=\event_0\cap\mathsf{Good}$. Then
%\[
%\Pr[\event]\geq 1-\epsilon_1-\epsilon_2-2/\lambda.
%\]
%%for any $\widehat{\vec s}\in\mathcal S^*$.
%%The first equation is due to the fact that $\vec s=\widehat{\vec s}$ implies $\event_a$.
%%Since $\Pr[\event_a]=(1-\frac{1}{n})^\pnum$, we have
%%$$\Pr[\event_1\land\event_a]=\Pr[\event_1|\event_a]\Pr[\event_a]= (1-\epsilon_1-\epsilon_2-\frac{1}{2^k})(1-\frac{1}{n})^\pnum\geq 1-\epsilon_1-\epsilon_2-\frac{\pnum+1}{n}.$$
%The public keys determine $\vec{s}_{\lambda}$.
%%$$H(\key|g,\vec z=\widehat{\vec z},\vec s=\widehat{\vec s}),\event)=H(\key|g,\vec z=\widehat{\vec z},\{h_i\}_{i\in[\pnum]}=\{\widehat h_i\}_{i\in[\pnum]},f=\widehat f,\event)$$
%%for any $(\widehat h_i,\widehat f)_{i\in[\pnum]}\in\mathcal F^*$.
%%which implies
%%$$
%%H(\key)-H(\key|g,\vec z=\widehat{\vec z},(h_i)_{i\in[\pnum]}=(\widehat h_i)_{i\in[\pnum]},f=\widehat f,\event)\leq \amem.
%%$$
%%Then one can easily see that we have
%%$$
%%H(\key)-H(\key|g,\vec z,(h_i)_{i\in[\pnum]},f,\event)\leq \amem.
%%$$
%Let $W$ denote all remaining public components other than $(g_\pnum,\vec{s}_{\lambda})$.
%Conditioned on $(g_\pnum,\vec z,\vec{s}_{\lambda},\event)$, the variable $W$ is independent of $\urs_{\vec{s}_{\lambda}}$ and hence of $\key=g_\pnum(\urs_{\vec{s}_{\lambda}})$. Thus, by data processing, we have
%%$$
%%H(\key|\vec z,(\pk_i)_{i\in[\pnum]},\event)=H(\key|g_\pnum,\vec z,\vec s,\event),
%%$$
%\[\begin{split}
% I(\key;(\vec z,(\pk_i)_{i\in[\pnum]})|\event)\leq I(\key;g_\pnum,\vec z,\vec{s}_{\lambda}|\event)\leq \theta.
%%=&H(\key)-H(\key|g_\pnum,\vec z,(h_i,g_i)_{i\in[\pnum]},\event).
%%=&H(\key)-\sum_{\widehat{\vec z}\in\mathcal Z^*,((\widehat h_i)_{i\in[\pnum]},f)\in\mathcal F^*}\Pr[\vec z=\widehat{\vec z},(h_i)_{i\in[\pnum]}=(\widehat h_i)_{i\in[\pnum]},f=\widehat f]\cdot \\
%%    &H(\key|g,\vec z=\widehat{\vec z},(h_i)_{i\in[\pnum]}=(\widehat h_i)_{i\in[\pnum]},\event)\\
%%=&\sum_{\widehat{\vec z}\in\mathcal Z^*,((\widehat h_i)_{i\in[\pnum]},f=\widehat f)\in\mathcal F^*}\Pr[\vec z=\widehat{\vec z},(h_i)_{i\in[\pnum]}=(\widehat h_i)_{i\in[\pnum]},f=\widehat f]\cdot \\
%%     &(H(\key)-H(\key|g,\vec z=\widehat{\vec z},(h_i)_{i\in[\pnum]}=(\widehat h_i)_{i\in[\pnum]},f=\widehat f,\event))\\
%     %=&H(\key)-H(\key|g,\vec z=\widehat{\vec z},(h_i)_{i\in[\pnum]}=(\widehat h_i)_{i\in[\pnum]},f=\widehat f,\event)\leq \amem,
%\end{split}\]
%Therefore, $\bsmke$ satisfies $(\amem,\epsilon_1+\epsilon_2+2/\lambda,\theta)$-security in the BSM.
%%completing this part of proof.
%
%Putting all the above together, \cref{thm:bsmke} immediately follows.
\end{proof}


%Due to the page limit, we refer the reader to \cref{sec:bsmke-proof} for the proof of \cref{thm:bsmke}.
%
%%Before proving \cref{thm:bsmke}, we give the following lemma serving as the core lemma of our construction.
%
%
%
%
%
%
%
%%\subsection{Construction}
%%In this section, we propose two constructions of fine-grained NIZ\pnum for circuit SAT. Both constructions run in $\aczerotwo$ and are secure against adversaries in $\aczero$.
%
%
%%The first construction combines a fine-grained OR-proof (instantiated as in \cref{sec:ornizk}) and the DVV encryption \cite{C:DegVaiVas16} (in a non black-box way). It supports statement circuits in $\aczerotwo$ and can also support statement circuits in $\aczero$ if we allow the prover run in $\aczero$, i.e., its functionality is fully-fledged. The proof size is $l\cdot O(\lambda^2)$, where $l$ denotes the number of all wires in the statement circuit.
%
%%In this section, we propose a fine-grained NIZ\pnum for $\aczero$ circuit SAT running in $\aczero$ and secure against adversaries in $\aczero$.
%
%
%%It supports statement circuits in $\aczerotwo$ and can also support statement circuits in $\aczero$ if we allow the prover run in $\aczero$, i.e., its functionality is fully-fledged. The proof size is $l\cdot O(\lambda^2)$, where $l$ denotes the number of all wires in the statement circuit.
%
%%The second construction is generically constructed from a new fine-grained sFHE scheme that we give later and fine-grained NIZ\pnums (instantiated as in \cref{sec:lnizk,sec:ornizk}). While it only support statements in $\aczerocm$ (even if we allow the prover run in $\aczero$), the proof size is independent with the witness size and only depend on the input size. Specifically, the size is $n\cdot O(\lambda^2)$ when the witness size is $n$.
%
% 
%%\subsection{Fully-Fledged Construction}
%%\subsection{The DVV Encryption} \label{sec:dvv}
%%In this section, we give the definition of generalized affine MAC, which generalizes the notion of standard affine MAC \cite{C:Bla\pnumilPan14} by using predicate encodings, and show how to construct it in the fine-grained setting under the assumption $\assump$.
%
%
%%\begin{lemma}
%%For $b_0,b_1,b_2\in\bin$, $b_2=\lnot(b_0\land b_1)$ iff 
%%$$0=1+b_0b_1-b_2$$
%%or
%%$$0=1+b_1+(b_0+1)b_1-b_2.$$
%%\end{lemma}
%
%%Before giving our construction, we prove the following theorem, which is necessary to show that our NIZ\pnum can be executed in $\aczero$.
%
%%By slightly changing the construction of $\{\fz\}$, we immediately achieve another family of circuits $\{\fo\}$ outputting the index of the first $1$-bits of its input.
%%\begin{figure}[h!]
%%\begin{center}
%%\framebox{
%%\small
%%  \begin{tabular}{l}
%%    \begin{minipage}[t]{0.95\textwidth}
%%    \underline{$\fo(b_1,…,b_n)$:}\\
%% For each $i\in[n]$, we compute $\hx_i =  \vec i \cdot b_i$ in parallel. 
%% 
%% For each $i\in[n]$, we compute $\vec y_i = \hx_i \cdot (1-\OR_{1\leq k \leq (1-i), 1\leq j \leq \adv}  (x_{k,j})) $.
%% 
%%Compute $\vec y_{i^*}=\OR_{1\leq i \leq n} (y_{i,1}) || … || \OR_{1\leq i \leq n} (y_{i,\adv})$.
%%
%%    \end{minipage} 
%%      \end{tabular}
%%      }
%%\end{center}
%%\caption{Definition of $\fo$. By $\vec i\in\bin^\adv$ we denote the bit-string representing the index $i$, where we assume the bit-representation of $n$ has $\adv$ bits. By $y_{i,j}$ we denote the $j$-th bit of $\vec y_i$.}
%%\label{fig:fz}
%%\end{figure}
%
%%\heading{Construction of Our NIZ\pnum.}
%%We now define the following languages
%%\[\lang=\{  x:\exists \vec w \in \bin^{\lambda-1}, \text{ s.t. }  x= \matE \vec w \}\]
%%%where $\matr M=\begin{pmatrix}\matE & \matr 0\\ \matr 0& \matE\end{pmatrix}$ for $\matE\in\ZeroSamp$.
%%and
%%\[\begin{split}\Lor_\lambda=\{ (\hx_i)_{i\in[\adv]}:&\exists \vec w \in \bin^{2\lambda}\ \text{ s.t. }\ \bigvee_{i\in[\adv]}\hx_i= \matr M' \vec w \
%%\\  \mbox{or}\ & \exists \vec w \in \bin^{(\adv+1)\cdot\lambda}\ \text{ s.t. }\ \hx_{(\adv+1)}=\matr M_2\vec w \}\end{split}\]
%%where 
%%$$\matr M'=\begin{pmatrix}\matE &  \matr 0\\ \matr 0& \matE\end{pmatrix}\in\bin^{2\lambda\times2(\lambda-1)}$$
%% and 
%% $$\matr M_2=\begin{pmatrix}\matE & \matr 0 & \cdots& \matr 0\\
%% \matr 0 & \ddots & \ddots& \vdots \\
%% \vdots& \ddots& \ddots &\matr 0\\
%%  \matr 0& \cdots & \matr 0 & \matE\end{pmatrix}\in\bin^{(\adv+1)\cdot\lambda\times(\adv+1)\cdot(\lambda-1)},$$
%%i.e., $\matr M'$ and $\matr M_2$ contain $\matE$'s in the main diagonal and $\matr 0$ in the other positions. One can see that $\{\lang\}$ and $\{\Lor_\lambda\}$ are supported by our NIZ\pnum for linear languages in \cref{sec:lnizk} and our OR-proof given in \cref{sec:ornizk} respectively.
%
%
%
%%Let $\{\lang^\aczero\}$ be any family of languages such that for all $\x\in\lang^\aczero$, we can run $\relation_\lambda^\aczero(\x,\cdot)$ in $\aczero$, where $\relation_\lambda^\aczero(\x,\cdot)$ is the associated relation. 
%%\footnote{We can assume that each $\relation^\aczero_\lambda(\x,\cdot)$ consists only of ${\AND}$ and $\OR$ gates, since by De Morgan Rules, we can move all $\NOT$ gates to just the inputs and the resulting circuit is still in $\aczero$. However, this may cause loss on efficiency.}
%%Without loss of generality, we assume that all the $\AND$ and $\OR$ gates have fan-in of some polynomial $\adv=\adv(\lambda)$.
%%Let $\LNIZ\pnum=\{\NIZ\pnumGen,\NIZ\pnumProve,\NIZ\pnumVer\}$ and $\ORNIZ\pnum=\{\ORGen,\ORProve,\allowbreak\ORVer\}$ be NIZ\pnums with simulators $\{\NIZ\pnumTGen,\NIZ\pnumSim\}$ and $\{\ORTGen,\allowbreak\ORSim\}$ for $\{\lang\}$ and $\{\Lor_\lambda\}$ respectively.
%%%\[\Lor_{\matE}=\{ \hx_0,\hx_1:\exists \vec w \in \bin^{2t}, \text{ s.t. } \hx_0= \matr M \vec w \lor \hx_1=\matr M \vec w\}\]
%%%where $\matr M=\begin{pmatrix}\matE & \matr 0\\ \matr 0& \matE\end{pmatrix}$ for $\matE\in\ZeroSamp$.
%%We give our NIZ\pnum for $\{\lang^\aczero_\lambda\}$ and its simulator in \cref{fig:con-acnizk,fig:con-acnizk-sim} respectively.
%
%\begin{proof}
%\heading{Complexity.}
%%<<<<<<< HEAD
%%First we note that each $\party_{i\in[\pnum-1]}$ (respectively, $\party_\pnum$) consumes $4\log k+q\cdot n/k$ (respectively, $2\log n+2\log k+q$) bits of memory to store bits in $\urs$ and at most two strongly 2-universal hash functions in $\mathcal H_n$ and one in $\mathcal H_{n/k}$ for comparison at any point in the protocol. %$\party_\pnum$ consumes $2\log n$ bits of memory to store $h_\pnum$ and $f$.
%%Memory used for $g$ is not counted since all $h_i$ and $f$ can be deleted before generating or receiving $g$, and the users can use the freed memory for it.
%%%The memory for storing the function in $\mathcal G_r$ is not counted since the users has deleted the function(s) in $\mathcal F_k$ from the memory at that time and can reuse the memory for storing the function in $\mathcal F_k$ to store the one in $\mathcal G_r$.
%%%Specifically, each $\party_{i\in[\pnum-1]}$ stores at most $q\cdot \frac{n}{n}$ bits in $\urs$ and two functions in $\mathcal F_k$, where each function in $\mathcal H_n$can be represented by only $2\log k$. $\party_\pnum stores$ $q$ bits in $\urs$, two functions in $\mathcal F_k$ and $\mathcal F_{\frac{n}{n}}$ respectively.
%%The communication between all parties only involves $(h_i)_{i\in[\pnum]}$, $f$, and $g$, which can be represented by $2\pnum\log k+2\log (n/k)+\log \lambda$ bits in total.
%%%
%%%%$q = \sqrt[\pnum]{\frac{2\lambda}{n}}\cdot k\cdot\frac{n}{n}=n \sqrt[\pnum]{\frac{2\lambda}{n}}$
%%%
%%%%Next we prove that the probability that $\party_i$ aborts for some $i\in[\pnum]$ is at most $\exp(-\lambda\delta^2/2)$.
%%%
%%%%Assuming that $h_i$ is generated as $h_i\getsr \mathcal H_u$ rather than $h_i\getsr \mathcal H$ for all $i\in[\pnum]$, and $\event_a$ be the event that $h_i\in\mathcal H$ for all $i\in[\pnum]$.  Then we have 
%%%%$$\Pr[\event_a]\geq (1-\frac{1}{2^k})^{\pnum}=\gamma.$$
%%%%Let $\set_i=(h_i(j))_{j\in\Z_q}$ for all $i\in[\pnum]$.
%%=======
%First we note that each $\party_{i\in[\pnum]}$ consumes $4\log n+\log\lambda+q$ bits of memory to store q bits in $\urs$, a strong universal hash function in $\mathcal G_{\lambda,r}$, and at most two strongly 2-universal hash functions in $\mathcal H_n$ for comparison at any point in the protocol. %$\party_\pnum$ consumes $2\log n$ bits of memory to store $h_\pnum$ and $f$.
%%Memory used for $g$ is not counted since all $h_i$ and $f$ can be deleted before generating or receiving $g$, and the users can use the freed memory for it.
%%The memory for storing the function in $\mathcal G_r$ is not counted since the users has deleted the function(s) in $\mathcal H_n$ from the memory at that time and can reuse the memory for storing the function in $\mathcal H_n$ to store the one in $\mathcal G_r$.
%%Specifically, each $\party_{i\in[\pnum-1]}$ stores at most $q\cdot \frac{n}{n}$ bits in $\urs$ and two functions in $\mathcal H_n$, where each function in $\mathcal H_n$can be represented by only $2\log n$. $\party_\pnum stores$ $q$ bits in $\urs$, two functions in $\mathcal H_n$ and $\mathcal F_{\frac{n}{n}}$ respectively.
%%The communication between all parties only involves $(h_i)_{i\in[\pnum]}$, $f$, and $g$, which can be represented by $2\pnum\log n+2\log (n/k)+\log \lambda$ bits in total.
%
%%$q = \sqrt[\pnum]{\frac{2\lambda}{n}}\cdot k\cdot\frac{n}{n}=n \sqrt[\pnum]{\frac{2\lambda}{n}}$
%
%%Next we prove that the probability that $\party_i$ aborts for some $i\in[\pnum]$ is at most $\exp(-\lambda\delta^2/2)$.
%
%%Assuming that $h_i$ is generated as $h_i\getsr \mathcal H_u$ rather than $h_i\getsr \mathcal H$ for all $i\in[\pnum]$, and $\event_a$ be the event that $h_i\in\mathcal H$ for all $i\in[\pnum]$.  Then we have 
%%$$\Pr[\event_a]\geq (1-\frac{1}{2^k})^{\pnum}=\gamma.$$
%%Let $\set_i=(h_i(j))_{j\in\Z_q}$ for all $i\in[\pnum]$.
%%>>>>>>> 55f330a20141b67cc25a1c99843653429dfe56a9
%\heading{Correctness.} To prove correctness, it is sufficient to prove that the probability that the protocol aborts is at most $\frac{2}{\lambda}$.
%
%Without loss of generality, we assume that $h_\pnum$ is randomly sampled from the strongly 2-universal hash function family $\mathcal F_n$ instead (see \cref{fig:con-hash}).
%Defining the indictor $Y_i$ as
%$Y_{i}=\begin{cases} 1 & x_i=h_\pnum(i)\in\cap_{l=1}^{\pnum-1}\set_l\\
%0 & \text{otherwise}
%\end{cases}$,
%%Defining the indictor $Y_i$ as $Y_i=1$ if 
%%$$x_i=h_\pnum(i)\in\cap_{l=1}^{\pnum-1}\set_l$$ and $Y_i=0$ otherwise,
%we have
%\[\begin{split}
%\Pr[Y_i=1]=&\Pr[ x_i\in\cap_{l=1}^{\pnum-1}\set_l]
%=\prod_{l=1}^{\pnum-1}\Pr[ x_i\in\set_l]=(\frac{q}{n})^{\pnum-1},
%\end{split}\]
%and for any $i\neq j$, we have
%\[\begin{split}
%&\Pr[Y_i=1\land Y_j=1]\\
%=&\Pr[x_i,x_j\in \cap_{l=1}^{\pnum-1}\set_l]\\
%=&\prod_{l=1}^{\pnum-1}\Pr[x_i,x_j\in \set_l]\ ( \because \mbox{independence of $\set_1,\cdots \set_{\pnum-1}$})\\
%=&\prod_{l=1}^{\pnum-1}(\Pr[x_i\in\set_l]\cdot\Pr[x_j\in \set_l]) \ ( \because \mbox{pairwise independence of $x_i$ and $x_j$})\\
%=&\Pr[x_i\in\cap_{l=1}^{\pnum-1}\set_l]\cdot\Pr[x_j\in \cap_{l=1}^{\pnum-1}\set_l]
%=\Pr[Y_i=1]\cdot \Pr[Y_j=1],
%\end{split}\]
%i.e., $Y_i$ and $Y_j$ are pairwise independent.
%
%Moreover, we have
%$$\expect[Y_i]=\Pr[Y_i=1]\cdot 1+\Pr[Y_i=0]\cdot 0=\Pr[Y_i=1]=(\frac{q}{n})^{\pnum-1},$$
%$$\expect[Y_i^2]=\Pr[Y_i^2=1]\cdot 1+\Pr[Y_i^2=0]\cdot 0=\Pr[Y_i^2=1]=\expect[Y_i],$$
%$$\var[Y_i]=\expect[Y_i^2]-\expect[Y_i]^2\geq\expect[Y_i]=(\frac{q}{n})^{\pnum-1}.$$
%
%%Hence the expected value of $|\sk_1|$ (which is equivalent to $|\sk_2|=\cdots=|\sk_\pnum|$) is
%%$$
%%\sum_{ x\in[n]}\Pr[\msb_\kk( x)\in\set_1\cap\set_2\cap\cdots\cap\set_\pnum]=n\cdot (\frac{q}{2^k})^\pnum=\lambda.
%%$$
%%Then the probability that $\party_i$ aborts for some $i$ is $\exp(-\lambda\delta^2/2)$ due to Chernoff bound.
%
%According to \cref{thm:cheb}, we have
%\[\begin{split}
%\Pr\Big[|Y_i\cdot q-\expect[Y_i]\cdot n|\geq \frac{\expect[Y_i]\cdot q}{2}\Big]\leq \frac{4}{q\expect[Y_i]},
%\end{split}\]
%$$
%\mbox{i.e.,}\ \Pr\Big[Y_i\cdot q\leq \frac{\expect[Y_i]\cdot q}{2}=\lambda\Big]\leq \frac{4}{q\expect[Y_i]}=\frac{2}{\lambda}.
%$$
%
%%i.e.,
%%\[\begin{split}
%%\Pr\Big[|Y_i\cdot n-2\lambda|\geq \lambda|\event_a\Big]\leq \frac{2}{\lambda\Pr[\event_a]}.
%%\end{split}\]
%%Then we have 
%%\[\begin{split}
%%\Pr\Big[Y_i\cdot n\geq \lambda|\event_a\Big]\leq \frac{2}{\lambda\cdot\gamma}.
%%\end{split}\]
%Since the number of indices shared by the parties when $h_\pnum\getsr\mathcal H_n$ must be more than that when $h_\pnum\getsr\mathcal F_n$, the probability the protocol aborts is at most $\frac{2}{\lambda}$, completing this part of proof.

%\heading{$(m,\epsilon_1+\epsilon_2),\amem)$-security.}
%Let $f:\bin^n\rightarrow\bin^\amem$ be any storage function applied by the adversary. %Let $\{ h_1(j)\}_{j\in\Z_q},\cdots,\{ h_\pnum(j)\}_{j\in\Z_q}$ be sequences of approximate pairwise independent random variables generated as $t_i^j=h_i(j)$ where $h_i\getsr \mathcal H_n$ for all $i\in[\pnum]$ and all $j\in\Z_q$.
%%Let $\{ h_\pnum(j)\}_{j\in\Z_q}$ be a sequence of $\kk$-approximate pairwise independent random variables with size $\nn$. 
%Let $\vec u=(u_i)_{i\in[q]}$ be the sequence such that $u_i= h_\pnum(i)$ if there are indices $i_1,\cdots,i_{\pnum-1}$ satisfying
%$ h_1(i_1)=\cdots= h_{\pnum-1}(i_{\pnum-1})= h_\pnum(i)$ and $ u_i=\bot$ otherwise, and let $\vec s$ be the sequence $\vec u$ restricted to $u_i\neq \bot$.
%We have the following lemma.
%The proof of this lemma is given in \cref{sec:pairwise}
%\begin{lemma}
%\label{lem:pairwise}
%$\vec s$ is generated by $\mathcal H_n$.
%%Then the sequence $(s_i)_{i\in[q]}$ restricted to those positions different from $\bot$ is pairwise independent.
%\end{lemma}

%\begin{proof}
%%Let $ t_i^j=h_i(j)$ for all $i\in[\pnum]$ and $j\in\Z_q$ and $( u_j)_{j\in\Z_q}$ be the sequence such that $ u_j= h_1(j)$ if there exist indices $j_2,\cdots,j_{\pnum}$ satisfying
%%$\msb_\kk( h_1(j))=\msb_\kk( t_i^{j_2})=\cdots=\msb_\kk( t_\pnum^{j_\pnum})$ and $ u_i=\bot$ otherwise.
%Let $\hx_1,\hx_2$ be any distinct elements in $GF(n)$.
%%To prove this lemma, it is sufficient to show that 
%%the sequence $( u_i)_{i\in[q]}$ restricted to $ u_i\neq\bot$ is $\kk$-approximately pairwise independent, i.e. $( u_i)_{i\in[q]}$ satisfies
%%$
%%\Pr[ u_i=\hx_1\land  u_j=\hx_2| u_i\neq \bot\land  u_j\neq\bot]=\frac{1}{k(k-1)}
%%$
%%for all $i,j\in\Z_q$. 
%Due to approximately pairwise independence, we have 
%$
%\Pr[ h_\pnum(i)=\hx_1\land  h_\pnum(j)=\hx_2]=\frac{1}{n(n-1)}
%$
%for all $i,j\in[q]$,
%and
%$
%\Pr[\exists i_l,j_l: h_l(i_l)=\hx_1\land h_l(j_l)=\hx_2]=\frac{q(q-1)}{n(n-1)}
%$
%for all $l\in[\pnum-1]$.
%Hence, for any $i,j\in\Z_q$, we have
%\[\begin{split}
%&\Pr[ u_i=\hx_1\land  u_j=\hx_2]\\
%=&\Pr[\exists (i_l,j_l)_{l\in[\pnum-1]}: 
% h_1(i_1)=\cdots= h_{\pnum-1}(i_{\pnum-1})=\hx_1= h_\pnum(i)\land\\
% &\tab h_1(j_1)=\cdots= h_{\pnum-1}(j_{\pnum-1})=\hx_2= h_\pnum(j)]\\
%=&\Pr[ h_\pnum(i)=\hx_1\land  h_1(j)=x_2]\cdot\prod_{l=1}^\pnum\Pr[\exists i_l,j_l: h_l(i_1)=\hx_1\land  h_l(j_l)=\hx_2]\\
%%=\Pr[&t_1^i=x_1\land h_1(j)=x_2]\cdot\Pr[\exists i_2,\cdots,i_\pnum,j_2,\cdots,j_\pnum: \\&
%%\msb_\kk(t_2^{i_2})=\cdots=\msb_\kk(t_\pnum^{i_\pnum})=\msb_\kk(x_1)\land\\
%%&\msb_\kk(t_2^{j_2})=\cdots=\msb_\kk(t_\pnum^{j_\pnum})=\msb_\kk(x_2)]\\
%=&\frac{1}{n(n-1)}\cdot \Big(\frac{q(q-1)}{n(n-1)}\Big)^{\pnum-1},
%\end{split}\]
%\[\begin{split}
%&\Pr[ u_i\neq\bot\land  u_j\neq \bot]
%\\
%=&\Pr[\exists (i_l,j_l)_{l\in[\pnum-1]}: h_1(i_1)=\cdots= h_{\pnum-1}(i_{\pnum-1})= h_\pnum(i)\land\\
%&\tab  h_1(j_1)=\cdots= h_{\pnum-1}(j_{\pnum-1})= h_\pnum(j)]
%\\
%=&\prod_{l=1}^\pnum\Pr[\exists i_l,j_l: h_l(i_l)=h_\pnum(i)\land  h_l(j_l)= h_\pnum(j)]
%=\Big(\frac{q(q-1)}{n(n-1)}\Big)^{\pnum-1},
%\end{split}\]
%
%$$
%\mbox{i.e.,} \Pr[ u_i=x_1\land  u_j=x_2| u_i\neq \bot\land  u_j\neq\bot]=\frac{\Pr[ u_i=\hx_1\land  u_j=\hx_2]}{\Pr[ u_i\neq\bot\land  u_j\neq \bot]}=\frac{1}{n(n-1)}.
%$$
%Hence, the distribution of $(u_i,u_j)$ is identical to that of $(h_\pnum(i),h_\pnum(j))$.
%Moreover, for each $l\in\set\backslash\{i,j\}$ we have $u_l=h_\pnum(l)$ where $h_\pnum$ can be determined by $u_i$ and $u_j$. Hence, the distribution of $(u_i)_{i\in\set}$ is identical to that of $(h_\pnum(i))_{i\in\set}$, completing the proof of \cref{lem:pairwise}.
%\end{proof}
%We refer the reader to \cref{sec:pairwiselem} for the proof of \cref{lem:pairwise}.

%\begin{proof}
%%Let $ t_i^j=h_i(j)$ for all $i\in[\pnum]$ and $j\in\Z_q$ and $( u_j)_{j\in\Z_q}$ be the sequence such that $ u_j= h_1(j)$ if there exist indices $j_2,\cdots,j_{\pnum}$ satisfying
%%$\msb_\kk( h_1(j))=\msb_\kk( t_i^{j_2})=\cdots=\msb_\kk( t_\pnum^{j_\pnum})$ and $ u_i=\bot$ otherwise.
%Let $\hx_1,\hx_2$ be any distinct elements in $GF(n)$.
%%To prove this lemma, it is sufficient to show that 
%%the sequence $( u_i)_{i\in[q]}$ restricted to $ u_i\neq\bot$ is $\kk$-approximately pairwise independent, i.e. $( u_i)_{i\in[q]}$ satisfies
%%$
%%\Pr[ u_i=\hx_1\land  u_j=\hx_2| u_i\neq \bot\land  u_j\neq\bot]=\frac{1}{k(k-1)}
%%$
%%for all $i,j\in\Z_q$. 
%Due to approximately pairwise independence, we have 
%$
%\Pr[ h_\pnum(i)=\hx_1\land  h_\pnum(j)=\hx_2]=\frac{1}{n(n-1)}
%$
%for all $i,j\in[q]$,
%and
%$
%\Pr[\exists i_l,j_l: h_l(i_l)=\hx_1\land h_l(j_l)=\hx_2]=\frac{q(q-1)}{n(n-1)}
%$
%for all $l\in[\pnum-1]$.
%Hence, for any $i,j\in\Z_q$, we have
%\[\begin{split}
%&\Pr[ u_i=\hx_1\land  u_j=\hx_2]\\
%=&\Pr[\exists (i_l,j_l)_{l\in[\pnum-1]}: 
% h_1(i_1)=\cdots= h_{\pnum-1}(i_{\pnum-1})=\hx_1= h_\pnum(i)\land\\
% &\tab h_1(j_1)=\cdots= h_{\pnum-1}(j_{\pnum-1})=\hx_2= h_\pnum(j)]\\
%=&\Pr[ h_\pnum(i)=\hx_1\land  h_1(j)=x_2]\cdot\prod_{l=1}^\pnum\Pr[\exists i_l,j_l: h_l(i_1)=\hx_1\land  h_l(j_l)=\hx_2]\\
%%=\Pr[&t_1^i=x_1\land h_1(j)=x_2]\cdot\Pr[\exists i_2,\cdots,i_\pnum,j_2,\cdots,j_\pnum: \\&
%%\msb_\kk(t_2^{i_2})=\cdots=\msb_\kk(t_\pnum^{i_\pnum})=\msb_\kk(x_1)\land\\
%%&\msb_\kk(t_2^{j_2})=\cdots=\msb_\kk(t_\pnum^{j_\pnum})=\msb_\kk(x_2)]\\
%=&\frac{1}{n(n-1)}\cdot \Big(\frac{q(q-1)}{n(n-1)}\Big)^{\pnum-1},
%\end{split}\]
%\[\begin{split}
%&\Pr[ u_i\neq\bot\land  u_j\neq \bot]
%\\
%=&\Pr[\exists (i_l,j_l)_{l\in[\pnum-1]}: h_1(i_1)=\cdots= h_{\pnum-1}(i_{\pnum-1})= h_\pnum(i)\land\\
%&\tab  h_1(j_1)=\cdots= h_{\pnum-1}(j_{\pnum-1})= h_\pnum(j)]
%\\
%=&\prod_{l=1}^\pnum\Pr[\exists i_l,j_l: h_l(i_l)=h_\pnum(i)\land  h_l(j_l)= h_\pnum(j)]
%=\Big(\frac{q(q-1)}{n(n-1)}\Big)^{\pnum-1},
%\end{split}\]
%
%$$
%\mbox{i.e.,} \Pr[ u_i=x_1\land  u_j=x_2| u_i\neq \bot\land  u_j\neq\bot]=\frac{\Pr[ u_i=\hx_1\land  u_j=\hx_2]}{\Pr[ u_i\neq\bot\land  u_j\neq \bot]}=\frac{1}{n(n-1)}.
%$$
%Hence, the distribution of $(u_i,u_j)$ is identical to that of $(h_\pnum(i),h_\pnum(j))$.
%Moreover, for each $l\in\set\backslash\{i,j\}$ we have $u_l=h_\pnum(l)$ where $h_\pnum$ can be determined by $u_i$ and $u_j$. Hence, the distribution of $(u_i)_{i\in\set}$ is identical to that of $(h_\pnum(i))_{i\in\set}$, completing the proof of \cref{lem:pairwise}.
%\end{proof}

%Due to \cref{lem:pairwise}, one can see that the distribution of $\set_1\cap\cdots \cap\set_\pnum$ is identical to that of $\{h_\pnum(i))\}_{i\in\set}$.
%Let $\mathcal S^*$, $\mathcal Z^*$, and $\mathcal F^*$ be the sets of all possible values of $\vec s$, $\vec z=\A(\urs)$, and $(\{h_i\}_{i\in[\pnum]},f)$ respectively conditioned on that $\event_a$ occurs.
%Combining the above lemma with \cref{thm:cm}, there exists an event $\event$ such that
%$\Pr[\event]\geq 1-\epsilon_1-\epsilon_2$ and
%\[\begin{split}
%I(\key;g_\pnum|\vec z,\vec s,\event))
%%&=I(\key;g|\vec z=\widehat{\vec z},\vec s=\widehat{\vec s},\event))\\
%=H(\key)-H(\key|g_\pnum,\vec z,\vec s,\event)\leq \theta
%\end{split}\]
%where $\vec z=f(\urs)$.
%%for any $\widehat{\vec s}\in\mathcal S^*$. 
%%The first equation is due to the fact that $\vec s=\widehat{\vec s}$ implies $\event_a$.
%%Since $\Pr[\event_a]=(1-\frac{1}{n})^\pnum$, we have
%%$$\Pr[\event_1\land\event_a]=\Pr[\event_1|\event_a]\Pr[\event_a]= (1-\epsilon_1-\epsilon_2-\frac{1}{2^k})(1-\frac{1}{n})^\pnum\geq 1-\epsilon_1-\epsilon_2-\frac{\pnum+1}{n}.$$
%Moreover, since $\vec s$ is determined by $\{h_i\}_{i\in[\pnum]}$ and
%%$$H(\key|g,\vec z=\widehat{\vec z},\vec s=\widehat{\vec s}),\event)=H(\key|g,\vec z=\widehat{\vec z},\{h_i\}_{i\in[\pnum]}=\{\widehat h_i\}_{i\in[\pnum]},f=\widehat f,\event)$$
%%for any $(\widehat h_i,\widehat f)_{i\in[\pnum]}\in\mathcal F^*$.
%%which implies 
%%$$
%%H(\key)-H(\key|g,\vec z=\widehat{\vec z},(h_i)_{i\in[\pnum]}=(\widehat h_i)_{i\in[\pnum]},f=\widehat f,\event)\leq \amem.
%%$$
%%Then one can easily see that we have
%%$$
%%H(\key)-H(\key|g,\vec z,(h_i)_{i\in[\pnum]},f,\event)\leq \amem.
%%$$
%the public keys $(\pk_i)_{i\in[\pnum]}$ between all parties only contain $\{h_i,g_i\}_{i\in[\pnum]}$, we have 
%\[\begin{split}
%I(\key;(\vec z,(\pk_i)_{i\in[\pnum]},\event)
%=&H(\key)-H(\key|g_\pnum,\vec z,(h_i,g_i)_{i\in[\pnum]},\event)\leq \theta.
%%=&H(\key)-\sum_{\widehat{\vec z}\in\mathcal Z^*,((\widehat h_i)_{i\in[\pnum]},f)\in\mathcal F^*}\Pr[\vec z=\widehat{\vec z},(h_i)_{i\in[\pnum]}=(\widehat h_i)_{i\in[\pnum]},f=\widehat f]\cdot \\
%%    &H(\key|g,\vec z=\widehat{\vec z},(h_i)_{i\in[\pnum]}=(\widehat h_i)_{i\in[\pnum]},\event)\\
%%=&\sum_{\widehat{\vec z}\in\mathcal Z^*,((\widehat h_i)_{i\in[\pnum]},f=\widehat f)\in\mathcal F^*}\Pr[\vec z=\widehat{\vec z},(h_i)_{i\in[\pnum]}=(\widehat h_i)_{i\in[\pnum]},f=\widehat f]\cdot \\
%%     &(H(\key)-H(\key|g,\vec z=\widehat{\vec z},(h_i)_{i\in[\pnum]}=(\widehat h_i)_{i\in[\pnum]},f=\widehat f,\event))\\
%     %=&H(\key)-H(\key|g,\vec z=\widehat{\vec z},(h_i)_{i\in[\pnum]}=(\widehat h_i)_{i\in[\pnum]},f=\widehat f,\event)\leq \amem,
%\end{split}\]
%%completing this part of proof.
%
%Putting all the above together, \cref{thm:bsmke} immediately follows.
%\end{proof}
%\heading{The special case where $k=n$.} Although we require that $q\leq\min\{k,n/k\}$, one can easily see that our proof also holds if set $k=n$ and require that $q\leq n$ instead. In this case, the strongly 2-universal hash function family we use is exactly the one in \cref{fig:con-hash}.

\heading{Example with asymptotic complexity.} Let $n=\lambda^{\pnum+1}$, $\amem=\lambda^{\pnum+1}/2$, $\epsilon_1=1/\lambda^\pnum$, $\epsilon_2=1/\sqrt{\lambda}$, and $\theta = \negl(\lambda)$. We have $r=O(\lambda)$. In this case, each $\party_{i\in[\pnum]}$ consumes about
$q=n\sqrt[\pnum]{2\lambda/n}=O(\lambda^\pnum)$
%$$q\cdot \frac{n}{n} = \sqrt[\pnum]{\frac{2\lambda}{n}}\cdot k\cdot\frac{n}{n}=n \sqrt[\pnum]{\frac{2\lambda}{n}}=O(\lambda^{4-\frac{1}{\pnum}})$$ bits of memory and $\party_\pnum$ consumes 
%$q=k\sqrt[\pnum]{\frac{2\lambda}{n}}=O(\lambda^{2-\frac{1}{\pnum}})$
bits of memory and the communication cost, i.e., the total size of the public keys, is $O(\lambda)$.
%If we set $k=n$ (i.e., all $h_i$ are sampled from $\mathcal F_n$) instead, then $\party_{i\in[\pnum]}$ consumes
%$q=n\sqrt[\pnum]{\frac{2\lambda}{n}}=O(\lambda^{4-\frac{3}{\pnum}})$
%bits of memory.

\heading{Concrete example.}
We now give a concrete example of the system based on the data given in \cite{C:CacMau97}.
Assume that all users have access to a $40$ Gbit/s broadcast channel used for $2\times 10^5$ seconds (about two days), we have $n = 8.6 \times 10^{15}$. Let $\amem=4.5 \times 10^{15}$. With $\theta = 10^{-20}$ and error probabilities $\epsilon_1=10^{-20}$ and $\epsilon_2=10^{-3}$, the parameters are $\delta = 0.476$, $\rho = 0.077$, $\lambda = 1.3 \times 10^7$, and $r = 5.0 \times 10^5$. Except with probability about $10^{-3}$, Eve knows less than $10^{-20}$ bits about the $61$ KByte session key.
%In order to share $\lambda$ bits,
To obtain this $r$-bit session key in a $3$-party protocol, each party consumes about $1.2\times 10^{13}$ bits of memory, and the communication consists of about $3.9\times 10^7$ bits in total. %In a $5$-party protocol, each party consumes about $1.7\times 10^{14}$ bits of memory in total.
%The public communication between all parties consists of about $235.3$ bits in total. 
%Setting $k=2.08\times 10^8$, to share $\lambda$ bits in a $3$-party protocol, $\party_1$ and $\party_2$ consumes about $4.3 \times 10^{15}$ bits of memory, while $\party_3$ only consumes $1.04\times 10^8$ bits of memory. The public communication between all parties consists of about $235.3$ bits in total. 
%If we set $k=n$, all parties consumes about $1.2\times 10^{13}$ bits of memory, and the public communication consists of about $336.4$ bits in total.

%\heading{Trade-off in the $2$-party setting.}
%The key exchange system in \cite{C:CacMau97} can be treated as a special case where we set $k=n$. In this case, to share a session key with the length $r = 5.0 \times 10^5$, each party consumes $3.3\times10^{11}$ as in \cite{C:CacMau97} and the total communication consists of $230.6$ bits. Our work demonstrates the trade-off that, e.g., if we set $k=1.04\times 10^8$, then one party consumes $4.3\times 10^{15}$ bits of memory, while another consumes only $5.2\times 10^7$ bits of memory, and the total communication consists of $178$ bits.


 
%Since $\lambda=\frac{n}{2}\cdot (\frac{q}{n})^\pnum$,
%in order to have $\lambda$ common indices, each $\party_{i\in[\pnum-1]}$ must store $n \sqrt[\pnum]{\frac{2\lambda}{n}}$
%%$$q\cdot \frac{n}{n} = \sqrt[\pnum]{\frac{2\lambda}{n}}\cdot k\cdot\frac{n}{n}=n \sqrt[\pnum]{\frac{2\lambda}{n}}$$
% bits in $\urs$ and $\party_\pnum$ must store $q=k\sqrt[\pnum]{\frac{2\lambda}{n}}$ bits in $\urs$. 
%In \cref{fig:comparison-mem}, we present the memory consumed by each party regarding different values of $k$ and different number of parties.
%Notice that the special case with $k=n$ and $\pnum=2$ is exactly the 2-party key agreement for public discussion.
%
%\begin{table}[htb]
%	\setlength\tabcolsep{0.3eM}
%	\begin{center}
%		\begin{minipage}[c]{\textwidth} \scriptsize
%			\begin{center}
%				\begin{tabular}{|l|l|l|l|l|l|}\hline
%					\multicolumn{1}{|l|}{} & \multicolumn{1}{c|}{$\pnum=2$} & \multicolumn{1}{c|}{$\pnum=3$} &\multicolumn{1}{c|}{$\pnum=5$} & \multicolumn{1}{c|}{$\pnum=10$} & \multicolumn{1}{c|}{$\pnum=100$}   \\
%					\hline
%					$k=2.6\times 10^{7}$ & $(8.6\times 10^{15},2.6\times 10^7)$ & $(8.6\times 10^{15},2.6\times 10^7)$ & $(8.6\times 10^{15},2.6\times 10^7)$& $(8.6\times 10^{15},2.6\times 10^7)$& $(8.6\times 10^{15},2.6\times 10^7)$   \\
%					$k=2.6\times 10^{11}$ & $(8.6\times 10^{13},2.6\times 10^9)$ & $(4.0\times 10^{14},1.2\times 10^{10})$ & $(1.4\times 10^{15},4.1\times 10^{10})$ & $(3.6\times 10^{15},1.0\times 10^{11})$  & $(7.8\times 10^{15},2.4\times 10^{11})$   \\
%					$k=8.6\times 10^{15}$ & $(4.7\times 10^{11},4.7\times 10^{11})$ & $(1.2\times 10^{13},1.2\times 10^{13})$ & $(1.7\times 10^{14},1.7\times 10^{14})$& $(1.2\times 10^{15},1.2\times 10^{15})$& $(7.1\times 10^{15},7.1\times 10^{15})$ \\
%					\hline
%				\end{tabular}
%				\vspace{10pt}
%				\caption{The amount of memory consumed by each party. $2\lambda\leq k\leq n$ and $\pnum$ denotes the number of parties. By $(x,y)$ we mean that each $\party_{i\in[\pnum-1]}$ consumes $x$ bits of memory and $\party_\pnum$ consumes $y$ bits of memory. Here we omit the amount of memory used to store the strongly 2-universal hash functions since it is quite small.
%				}
%				\label{fig:comparison-mem}
%			\end{center}
%		\end{minipage}
%	\end{center}
%%		\vspace*{-1cm}
%\end{table}
%
%The public communication for the selected indices consists of $2(\pnum-1) \log n+\log n+\log r$ bits in total and .

%remark on applications
\heading{Extensions.} Similar to our constructions in previous sections, our construction in the BSM can also be converted into a BPRF and a multi-party IB-NIKE. We refer the reader to \cref{sec:ibnike} for the details.
%We refer the reader to the full paper for the details.

%Parts of Jiaxin Pan's work was supported by the Research Council of Norway under Project No. 324235.


\subsection{Limitations of Multi-Party NIKEs in the BSM}
\label{sec:negative}
In this section, we show that for an NIKE treating any constant number of parties satisfying (restricted) extendability and local key independence, the entropy of the shared key conditioned on the public keys is bounded by roughly $\umem^\pnum/\amem^{\pnum-1}$, where $\umem$ and $\amem$  are the memory bounds for users and adversaries respectively. This suggests that $\umem$ must be at least about $\sqrt[\pnum]{\lambda\cdot\amem^{\pnum-1}}$ for sharing $\lambda$ bits, and thus our multi-party NIKE in the BSM, which satisfies local key independence as we show later, achieves optimality up to constant factors within this natural class.

Before giving our negative result, we recall several lemmata and theorems below.
\begin{lemma}
\label{fact}
For two random variables $x, y$, we have
$H(x) \geq H(x\mid y)$.
\end{lemma}

\begin{theorem}[Data processing inequality]
\label{prop:datap}
For two random variables $x, y$ and deterministic functions $f,g$, we have
\[
H(f(x))\leq H(x)
\ \mbox{and}\
I(f(x);g(y))\leq I(x;y).
\]
The same inequalities hold conditioned on any additional random variable.
\end{theorem}

\begin{lemma}[Lemma 1 in \cite{EC:DziMau04}]
\label{lem:DM04}
Let $y$, $z$, $(z_i)_{i\in[n]}$ be a sequence of random variables such that conditioned on $y$, the variables $z, z_1,\cdots,z_n$ are independent and identically distributed, i.e.,
%$$
%\Pr[z=\widehat{z}\land z_1=\widehat z_1\land z_n=\widehat z_n | y=\widehat y]=\Pr[z=\widehat z|y=\widehat y]\prod_{i=1}^n\Pr[z=\widehat z_i|y=\widehat y]
%$$
$$
\Pr[z=\widehat{z}\land z_1=\widehat z_1\land \cdots \land z_n=\widehat z_n \mid y=\widehat y]=\Pr[z=\widehat z\mid y=\widehat y]\prod_{i=1}^n\Pr[z=\widehat z_i\mid y=\widehat y]
$$
for all $y,z,(z_i)_{i\in[n]}$. Then
$
I(y;z\mid z_1,\cdots,z_n)\leq {H(y)}/({n+1}).
$
\end{lemma}

\begin{theorem}[Theorem 1 in \cite{Maurer93}]
\label{thm:M93}
Let $x,y,z,c_0,c,s,s'$ be random variables such that
\[
H(s\mid x,c_0,c)=0,\quad
H(s'\mid y,c_0,c)=0,\quad
\Pr[s\neq s']=0,
\]
and
\[
I(x;y\mid z,c_0,c)\leq I(x;y\mid z,c_0).
\]
Then
\[
H(s)\leq I(x;y\mid z,c_0)+I(s;z,c_0,c).
\]
\end{theorem}

\heading{Our negative result.} Below we give our negative result showing limitations of multi-party NIKE in the BSM. 
%By $p(\epsilon)$ we denote  $H_b(\epsilon)+\epsilon\log_2(|\set|-1)$ where $|\set|$ is the space of local key. One can see that $p(\epsilon)$ is equal to $0$ when $\epsilon=0$ as in our multi-party NIKE in the BSM.
\begin{theorem}
\label{thm:negative}
Let $\mathcal C_1$ be the class of streaming word-RAMs with memory bounds $\umem$.
%Let  $\bsmke=\{\gen,\share\}$ be a $\mathcal C_1$-$\pnum$-NIKE for any constant $\pnum\geq 2$ in the BSM with URS length $n$ and satisfying $\epsilon$-restricted extendability and $\mu$-local key independence, where the associated combining and extracting algorithm families are $\{\combine,\extract\}$. Then there exists an event $\event$ such that $\Pr[\event]\geq 1-\epsilon$ and conditioned on $\event$ we have %Fixing $\urs\in\bin^n$, there exists some event $\event$ such that $\Pr[\event]\geq 1-\delta$ and
Let $\mathsf{NIKE}=\{\gen,\share\}$ be a $\mathcal C_1$-$\pnum$-NIKE for any constant $\pnum\geq 2$ in the BSM with URS length $n$ and satisfying $0$-restricted extendability and $((\mu_i)_{i\in[\pnum]},\lfloor\amem/\umem\rfloor)$-local key independence for any $\amem\geq \umem$, where the associated combining and extracting algorithm families are $\{\combine,\extract\}$.\footnote{We use $t=\lfloor \amem/\umem\rfloor$. This only affects constant factors; for readability one may assume that $\amem/\umem$ is an integer.} Then
$$
H(\key\mid(\pk_i)_{i\in[\pnum]})
\leq
\umem^\pnum/\amem^{\pnum-1}
+
2\sum_{i=2}^{\pnum}(\umem/\amem)^{\pnum-i}\mu_i,
$$
where $\urs\getsr\bin^n$, $(\pk_i,\sk_i)\getsr\gen(\urs)$ for all $i\in[\pnum]$, and $\key=\share(1,\sk_1,\allowbreak(\pk_i)_{i\in[\pnum]})$.
\end{theorem}

\begin{proof}
Let $\localkey_{[i]}=\combine(\localkey_{[i-1]},\pk_i)$ for all $i>1$, where
$\localkey_{[1]}=\sk_1$. Let $M_i=(\pk_i,\sk_i)$ and
$\localkey_{[i]}'=\combine(\sk_i,(\pk_j)_{i>j\geq 1})$ for all $i\in[\pnum]$.
Let $t=\lfloor\amem/\umem\rfloor$, sample
$M_i'=(\pk_i',\sk_i')\getsr\gen(\urs)$ for all $i\in[t]$, and set
$M_E=(M_i')_{i\in[t]}$.

Fix any $i>1$. We apply \cref{thm:M93} with
\[
x=\localkey_{[i-1]},
y=M_i,
z=M_E,
c_0=(\pk_j)_{j\in[i-1]},
c=\pk_i,s=\localkey_{[i]},s'=\localkey_{[i]}'.
\]
By $0$-restricted extendability and determinism of $\combine$, we have
\[
\begin{gathered}
H(\localkey_{[i]}\mid \localkey_{[i-1]},(\pk_j)_{j\in[i-1]},\pk_i)=0,\\
H(\localkey_{[i]}'\mid M_i,(\pk_j)_{j\in[i-1]},\pk_i)=0,\\
\Pr[\localkey_{[i]}\neq \localkey_{[i]}']=0.
\end{gathered}
\]
Moreover, since $H(\pk_i\mid M_i)=0$, we have
\[
\begin{aligned}
&I(\localkey_{[i-1]};M_i\mid M_E,(\pk_j)_{j\in[i-1]})\\
=&
I(\localkey_{[i-1]};\pk_i,M_i\mid M_E,(\pk_j)_{j\in[i-1]})\\
=&
I(\localkey_{[i-1]};\pk_i\mid M_E,(\pk_j)_{j\in[i-1]})+
I(\localkey_{[i-1]};M_i\mid M_E,(\pk_j)_{j\in[i-1]},\pk_i)\\
&\geq
I(\localkey_{[i-1]};M_i\mid M_E,(\pk_j)_{j\in[i-1]},\pk_i).
\end{aligned}
\]
Therefore, by \cref{thm:M93} and local key independence, we have
\[
H(\localkey_{[i]})
\leq
I(\localkey_{[i-1]};M_i\mid M_E,(\pk_j)_{j\in[i-1]})+\mu_i.
\]

Next, by the second condition in \cref{def:localkeyind}, we have
\[
I(\localkey_{[i-1]};M_i\mid M_E,(\pk_j)_{j\in[i-1]})
\leq
I(\localkey_{[i-1]};M_i\mid M_E,(\pk_j)_{j\in[i-1]},\urs)+\mu_i.
\]
Conditioned on $\urs$ and $(\pk_j)_{j\in[i-1]}$, the variables
$M_i,M_1',\ldots,M_t'$ are independent copies generated in the same way.
Hence, by \cref{lem:DM04}, we have
\[
I(\localkey_{[i-1]};M_i\mid M_E,(\pk_j)_{j\in[i-1]},\urs)
\leq
\frac{H(\localkey_{[i-1]}\mid \urs,(\pk_j)_{j\in[i-1]})}{t+1}
\leq
\frac{H(\localkey_{[i-1]})}{t+1}.
\]
Since $t=\lfloor\amem/\umem\rfloor$, we have $1/(t+1)\leq\umem/\amem$.
Combining the above inequalities gives
\[
H(\localkey_{[i]})
\leq
(\umem/\amem)H(\localkey_{[i-1]})+2\mu_i.
\]

We now iterate this recurrence. Since
$H(\localkey_{[1]})\leq H(M_1)\leq\umem$, we have
\[
H(\localkey_{[2]})
\leq
(\umem/\amem)H(\localkey_{[1]})+2\mu_2
\leq
\umem^2/\amem+2\mu_2.
\]
Assume that, for some $i\geq3$,
\[
H(\localkey_{[i-1]})
\leq
\umem^{i-1}/\amem^{i-2}
+
2\sum_{\ell=2}^{i-1}(\umem/\amem)^{i-1-\ell}\mu_\ell.
\]
Then we have
\[
\begin{aligned}
H(\localkey_{[i]})
\leq&
(\umem/\amem)H(\localkey_{[i-1]})+2\mu_i
\leq
\umem^i/\amem^{i-1}
+
2\sum_{\ell=2}^{i}(\umem/\amem)^{i-\ell}\mu_\ell.
\end{aligned}
\]
Thus,
\[
H(\localkey_{[\pnum]})
\leq
\umem^\pnum/\amem^{\pnum-1}
+
2\sum_{i=2}^{\pnum}(\umem/\amem)^{\pnum-i}\mu_i.
\]

Finally, conditioning on the public keys,
\[
\begin{aligned}
H(\key\mid(\pk_i)_{i\in[\pnum]})
=&
H(\extract(\localkey_{[\pnum]},\pk_\pnum)\mid(\pk_i)_{i\in[\pnum]})
\leq
H(\localkey_{[\pnum]}\mid(\pk_i)_{i\in[\pnum]})\\
\leq&
H(\localkey_{[\pnum]})
\leq
\umem^\pnum/\amem^{\pnum-1}
+
2\sum_{i=2}^{\pnum}(\umem/\amem)^{\pnum-i}\mu_i.
\end{aligned}
\]
This completes the proof of \cref{thm:negative}.
\end{proof}



\paragraph{\bf Optimality of $\bsmke$.}
%Let $\pnum$, $n$, $\lambda$, and $q$ be the parameters defined in \cref{sec:bsm-nike}. As in mentioned in \cref{sec:extendable}, our construction $\bsmke$  in \cref{fig:con-bsmke} satisfies $2/\lambda$-extendability.
%To show the optimality, we now only have to show the local key independence of  $\bsmke$.
Let $\pnum$, $n$, $\lambda$, and $q$ be as in \cref{sec:bsm-nike}. By \cref{sec:extendable}, $\bsmke$ satisfies $0$-extendability, while aborts are covered by the $2/\lambda$-correctness bound in \cref{thm:bsmke}. Thus, it remains to verify local key independence. The next lemma gives the required bounds and shows that our instantiation matches the lower bound up to constant factors.
%Let $\pnum$, $n$, $\lambda$, and $q$ be the parameters defined in \cref{sec:bsm-nike}. As mentioned in \cref{sec:extendable}, our construction $\bsmke$  in \cref{fig:con-bsmke} satisfies $0$-extendability.
%Its abort probability is accounted for separately by the $2/\lambda$-correctness guarantee in \cref{thm:bsmke}.
%It remains to verify the local key independence parameters for BSMKE. The following lemma shows that BSMKE satisfies the bounds needed to match the lower bound, and hence that our instantiation is optimal up to constant factors.
%\begin{lemma}
%\label{lem:localkeyind}
%$\bsmke$ (see \cref{fig:con-bsmke}) satisfies $\pnum (1+(q-\pnum)/2n)$-local key independence.
%\end{lemma}
\begin{lemma}
\label{lem:localkeyind}
For every $t\in\N$, $\bsmke$ satisfies $((\mu_i)_{i\in[\pnum]},t)$-local key independence in the BSM with
\[
\mu_i=O(t\cdot q^{i+1}/n^i+q^i/n^{i-1}+\log n) \ \mbox{for all $i\in[\pnum]$}.
\]
\end{lemma}

\begin{proof}
For notational simplicity, let $\set_i$ be the index set stored by the $i$-th user, and let
$\set_{[i]}=\set_1\cap\cdots\cap\set_i$, so that $\localkey_{[i]}=\urs_{\set_{[i]}}$.
We suppress the compact public hash descriptions defining these sets; their size is $O(\log n)$ and is absorbed into the $O(\log n)$ term.

Fix $i\in[\pnum]$, and let $\set'_j$ be the index set stored in $\sk'_j$. The public keys reveal the length of the variable-length string $\urs_{\set_{[i]}}$, which costs at most $O(\log n)$ bits. Once the public keys fix the index sets, the only correlation between $\urs_{\set_{[i]}}$ and $\urs_{\set'_1\cup\cdots\cup\set'_t}$ comes from their overlapping positions. Hence
\[
\begin{aligned}
&I\left(
\localkey_{[i]};
(\pk_j)_{j\in[i]},(\pk'_j,\sk'_j)_{j\in[t]}
\right)\\
=&
I\left(
\urs_{\set_{[i]}};
(\pk_j)_{j\in[i]},(\pk'_j)_{j\in[t]}
\right)
+
I\left(
\urs_{\set_{[i]}};
(\sk'_j)_{j\in[t]}
\mid
(\pk_j)_{j\in[i]},(\pk'_j)_{j\in[t]}
\right)\\
\leq&
O(\log n)
+
I\left(
\urs_{\set_{[i]}};
\urs_{\set'_1\cup\cdots\cup\set'_t}
\mid
(\pk_j)_{j\in[i]},(\pk'_j)_{j\in[t]}
\right)\\
\leq&
O(\log n)
+
\expect\left[\left|\set_{[i]}\cap(\set'_1\cup\cdots\cup\set'_t)\right|\right]\\
\leq&
O(\log n)+\sum_{j=1}^t \expect[|\set_{[i]}\cap\set'_j|]\\
=&
O(\log n)+t\sum_{a\in GF(n)}\Pr[a\in\set_{[i]}\cap\set'_1]\\
=&
O(\log n)+t\cdot n(q/n)^{i+1}\\
=&
O(t\cdot q^{i+1}/n^i+\log n).
\end{aligned}
\]
This proves the first condition.

It remains to verify the second condition. For $i>1$, conditioned on $\urs$ and $(\pk_j)_{j\in[i-1]}$, the value $\localkey_{[i-1]}=\urs_{\set_{[i-1]}}$ is fixed. Therefore,
\[
I(\localkey_{[i-1]};(\pk_i,\sk_i)\mid(\pk_j)_{j\in[i-1]},(\pk'_j,\sk'_j)_{j\in[t]},\urs)=0.
\]
Without conditioning on $\urs$, the new key pair $(\pk_i,\sk_i)$ can reveal information about $\localkey_{[i-1]}$ only through the overlap
$\set_{[i-1]}\cap\set_i=\set_{[i]}$. Hence
\[
\begin{aligned}
&I(\localkey_{[i-1]};(\pk_i,\sk_i)\mid(\pk_j)_{j\in[i-1]},(\pk'_j,\sk'_j)_{j\in[t]})
\leq
\expect[|\set_{[i]}|]
=
\sum_{a\in GF(n)}\Pr[a\in\set_{[i]}]
=
n(q/n)^i
=
q^i/n^{i-1}.
\end{aligned}
\]
Thus we have
\[
\begin{aligned}
&I(\localkey_{[i-1]};(\pk_i,\sk_i)\mid(\pk_j)_{j\in[i-1]},(\pk'_j,\sk'_j)_{j\in[t]})\\
\leq&
I(\localkey_{[i-1]};(\pk_i,\sk_i)\mid(\pk_j)_{j\in[i-1]},(\pk'_j,\sk'_j)_{j\in[t]},\urs)
+
q^i/n^{i-1}.
\end{aligned}
\]
Since
$
\mu_i=O(t\cdot q^{i+1}/n^i+q^i/n^{i-1}+\log n),
$
both conditions in \cref{def:localkeyind} are satisfied.
\end{proof}

\heading{\bf Remark.}
For the usual parameter choice $n=\Theta(\amem)$, $\umem=\Theta(q)$, and $t=\lfloor\amem/\umem\rfloor$, \cref{lem:localkeyind} gives
\[
t\cdot q^{i+1}/n^i
=
O((\amem/\umem)\cdot \umem^{i+1}/\amem^i)
=
O(\umem^i/\amem^{i-1})
\ 
\mbox{and}
\
q^i/n^{i-1}
=
O(\umem^i/\amem^{i-1}).
\]
Hence, we have
$
\mu_i=
O(\umem^i/\amem^{i-1}+\log n).
$
Plugging this into \cref{thm:negative} yields
\[
H(\key\mid(\pk_i)_{i\in[\pnum]})
\leq
O(\umem^\pnum/\amem^{\pnum-1}+\log n).
\]
When $\log n=o(\lambda)$, any scheme in this class with
$H(\key\mid(\pk_i)_{i\in[\pnum]})=\Omega(\lambda)$ must satisfy
\[
\umem=\Omega\left(\sqrt[\pnum]{\lambda\cdot\amem^{\pnum-1}}\right).
\]
Thus, the BSM construction is optimal up to constant factors within this class.
%When $n=\lambda^{\pnum+1}$, $\amem=\lambda^{\pnum+1}/2$, and $q=n\sqrt[\pnum]{\frac{2\lambda}{n}}=O(\lambda^\pnum)$ as in the example with asymptotic complexity we mentioned before, $\bsmke$ satisfies O(2).
%
%When, e.g., $n=\lambda^{\pnum+1}$, $\amem=\lambda^{\pnum+1}/2$, and $$. We have . In this case, each $\party_{i\in[\pnum]}$ consumes about
%$q=n\sqrt[\pnum]{\frac{2\lambda}{n}}=O(\lambda^\pnum)$
%%$$q\cdot \frac{n}{n} = \sqrt[\pnum]{\frac{2\lambda}{n}}\cdot k\cdot\frac{n}{n}=n \sqrt[\pnum]{\frac{2\lambda}{n}}=O(\lambda^{4-\frac{1}{\pnum}})$$ bits of memory and $\party_\pnum$ consumes 
%%$q=k\sqrt[\pnum]{\frac{2\lambda}{n}}=O(\lambda^{2-\frac{1}{\pnum}})$
%bits of memory and the communication cost, i.e., the total size of the public keys, is $O(\lambda)$.


\subsection*{\bf Acknowledgements}
Part of this work was supported by the National Natural Science Foundation of China under Grant Numbers 62472073 and 62272091.
